\documentclass[a4paper,10pt,landscape]{article}
\usepackage[left=5mm,right=5mm,top=7mm,bottom=7mm,includehead,headheight=11pt,headsep=3mm]{geometry}
\usepackage{multicol}
\usepackage{listings}
\usepackage{titlesec}
\usepackage{enumitem}
\usepackage{fancyhdr}
\usepackage[T1]{fontenc}
\usepackage{courier}
\usepackage{helvet}
\usepackage[expansion=false]{microtype}
\usepackage{amsmath}
\usepackage{amssymb}
\usepackage{xcolor}
\usepackage{hyperref}
\hypersetup{hidelinks}

% --- Page Style ---
\pagestyle{fancy}
\fancyhf{}
\renewcommand{\headrulewidth}{0.3pt}
\fancyhead[L]{\footnotesize\sffamily\bfseries Team: JUST BitWarriors / IUPC 2026}
\fancyhead[C]{\footnotesize\sffamily\bfseries TEAM REFERENCE NOTEBOOK 2026}
\fancyhead[R]{\footnotesize\sffamily\bfseries Page \thepage}

% --- Custom Layout Macros ---
\newcommand{\catsection}[1]{%
  \vspace{3pt}%
  \phantomsection\addcontentsline{toc}{section}{#1}%
  \noindent\colorbox{black!85}{\parbox{\dimexpr\linewidth-2\fboxsep\relax}{\color{white}\bfseries\sffamily\fontsize{8.5pt}{9.5pt}\selectfont\centering #1}}%
  \vspace{2pt}%
}
\newcommand{\algotitle}[1]{%
  \vspace{2pt}%
  \noindent{\sffamily\bfseries\fontsize{7.5pt}{8.5pt}\selectfont #1}%
  \vspace{0.5pt}%
  \hrule height 0.3pt%
  \vspace{1.5pt}%
}
\newcommand{\algodesc}[1]{%
  {\noindent\itshape\fontsize{6.5pt}{7.5pt}\selectfont #1\par}\vspace{1pt}%
}

% --- Code Listings Style ---
\lstset{
  language=C++,
  basicstyle=\ttfamily\fontsize{6.5pt}{7.4pt}\selectfont,
  keywordstyle=\bfseries\color{blue!80!black},
  commentstyle=\itshape\color{gray!75!black},
  stringstyle=\color{teal!70!black},
  frame=none,
  breaklines=true,
  breakatwhitespace=false,
  tabsize=2,
  showstringspaces=false,
  columns=fullflexible,
  keepspaces=true,
  aboveskip=1pt,
  belowskip=1.5pt,
  xleftmargin=0pt,
  xrightmargin=0pt,
  morekeywords={int32_t,long,int,auto,vector,pair,set,map,string,queue,stack,priority_queue,function,array,deque,bitset,multiset,unordered_set,unordered_map,uint64_t,size_t,struct,class,typename,template,__int128,__int128_t,mt19937,mt19937_64,ld},
  literate={->}{{$\rightarrow$}}1 {<=}{{$\le$}}1 {>=}{{$\ge$}}1 {!=}{{$\ne$}}1
}

\setcounter{tocdepth}{2}
\renewcommand{\contentsname}{\centering\Large\sffamily\bfseries TABLE OF CONTENTS}

\begin{document}

% =========================================================================
% PAGE 1: TITLE & TABLE OF CONTENTS
% =========================================================================
\thispagestyle{fancy}
\begin{center}
  {\LARGE\sffamily\bfseries IUPC / ICPC TEAM REFERENCE NOTEBOOK 2026}\\[2pt]
  {\small\sffamily Jashore University of Science and Technology \quad $\vert$ \quad Competitive Programming Reference Document}
\end{center}
\vspace{-4pt}
\hrule height 0.5pt
\vspace{4pt}

\begin{multicols*}{3}
\setlength{\columnsep}{6pt}
\setlength{\columnseprule}{0.2pt}
\tableofcontents
\end{multicols*}

\newpage

% =========================================================================
% CHEATSHEET & CODE
% =========================================================================
\begin{multicols*}{3}
\setlength{\columnsep}{6pt}
\setlength{\columnseprule}{0.2pt}
\raggedcolumns


\catsection{1. Setup, Fast I/O \& Debugging}

\algotitle{1.1 Contest Boilerplate (Marwan's Setup)}
\algodesc{Standard fast competitive programming template with 64-bit int and helper macros.}
\begin{lstlisting}
#include <bits/stdc++.h>
using namespace std;
#define int long long
#define M 1000000007
#define N 1000005
#define INF 1e17
#define endl "\n"
#define all(v) (v).begin(), (v).end()
#define yes cout << "Yes" << endl
#define no cout << "No" << endl
#define minus cout << "-1" << endl
#define zero cout << "0" << endl
#define make_unique(x) sort(all((x))); (x).erase(unique(all((x))), (x).end())

void marwan() {
  // Solve logic
}

int32_t main() {
  ios_base::sync_with_stdio(false);
  cin.tie(NULL);
  int t = 1;
  cin >> t;
  while (t--) marwan();
  return 0;
}
\end{lstlisting}

\algotitle{1.2 GCC Optimization Pragmas}
\algodesc{Include at top of file for heavy CPU tasks. May disable vectorization on some platforms.}
\begin{lstlisting}
#pragma GCC optimize("O3,unroll-loops")
#pragma GCC target("avx2,bmi,bmi2,lzcnt,popcnt")
\end{lstlisting}

\algotitle{1.3 Sublime Text Build Systems}
\algodesc{Automated compilation and local file redirection.}
\begin{lstlisting}
// Linux:
{
  "shell": true,
  "cmd": "g++ -std=c++17 '${file}' -o '${file_path}/${file_base_name}' && '${file_path}/${file_base_name}' < input.txt > output.txt",
  "working_dir": "${file_path}",
  "selector": "source.cpp"
}
// Windows:
{
  "cmd": ["g++.exe", "-std=c++17", "${file}", "-o", "${file_base_name}.exe", "&&", "${file_base_name}.exe<input.txt>output.txt"],
  "shell": true,
  "working_dir": "$file_path",
  "selector": "source.cpp"
}
\end{lstlisting}

\algotitle{1.4 Fast I/O (fread based)}
\algodesc{Reads integers $\sim 5\times$ faster than cin. Do NOT mix with cin/scanf.}
\begin{lstlisting}
inline char gc() {
  static char buf[1 << 16];
  static size_t bc = 0, be = 0;
  if (bc >= be) {
    buf[0] = 0; bc = 0;
    be = fread(buf, 1, sizeof(buf), stdin);
  }
  return buf[bc++];
}
int readInt() {
  int a, c;
  while ((a = gc()) < 40) if (a == 0) return 0;
  bool minus = (a == '-');
  if (minus) a = gc();
  int res = a - '0';
  while ((c = gc()) >= '0' && c <= '9') res = res * 10 + (c - '0');
  return minus ? -res : res;
}
\end{lstlisting}

\algotitle{1.5 Universal Debug Macro}
\algodesc{Prints variable names and container contents (vector, set, map, pair) to stderr.}
\begin{lstlisting}
#ifndef ONLINE_JUDGE
#define dbg(...) cerr << __LINE__ << ": ", __dbg(#__VA_ARGS__, __VA_ARGS__), cerr << "\n"
template<class T> void __p(T x) { cerr << x; }
template<class T, class U> void __p(pair<T, U> x) { cerr << "(" << x.first << "," << x.second << ")"; }
template<class T> void __p(vector<T> x) { cerr << "{"; int _=0; for (auto y : x) { if (_++) cerr << ","; __p(y); } cerr << "}"; }
template<class T> void __p(set<T> x) { cerr << "{"; int _=0; for (auto y : x) { if (_++) cerr << ","; __p(y); } cerr << "}"; }
template<class T, class U> void __p(map<T, U> x) { cerr << "{"; int _=0; for (auto [k, v] : x) { if (_++) cerr << ","; __p(k); cerr << ":"; __p(v); } cerr << "}"; }
void __dbg(const char* s) {}
template<class T, class... U> void __dbg(const char* s, T x, U... t) {
  int i = 0; while (s[i] && s[i] != ',') i++;
  cerr.write(s, i) << "="; __p(x);
  if constexpr (sizeof...(t)) cerr << " || ", __dbg(s + i + 1, t...);
}
#else
#define dbg(...)
#endif
\end{lstlisting}

\algotitle{1.6 GP Hash Table \& Anti-Hash Test}
\algodesc{High-performance hash table immune to worst-case anti-hash tests.}
\begin{lstlisting}
#include <ext/pb_ds/assoc_container.hpp>
using namespace __gnu_pbds;

struct custom_hash {
  static uint64_t splitmix64(uint64_t x) {
    x += 0x9e3779b97f4a7c15;
    x = (x ^ (x >> 30)) * 0xbf58476d1ce4e5b9;
    x = (x ^ (x >> 27)) * 0x94d049bb133111eb;
    return x ^ (x >> 31);
  }
  size_t operator()(uint64_t x) const {
    static const uint64_t FIXED_RANDOM = chrono::steady_clock::now().time_since_epoch().count();
    return splitmix64(x + FIXED_RANDOM);
  }
  size_t operator()(pair<int, int> p) const noexcept {
    return (*this)(p.first) ^ ((*this)(p.second) << 1);
  }
};
template<typename K, typename V>
using safe_map = gp_hash_table<K, V, custom_hash>;
template<typename K>
using safe_set = gp_hash_table<K, null_type, custom_hash>;
\end{lstlisting}

\algotitle{1.7 Basic Node Structures (Linked List \& Tree)}
\algodesc{Snippet definitions for pointers, linked lists, and tree nodes.}
\begin{lstlisting}
class Node {
public:
  int data;
  Node* next;
  Node(int d) : data(d), next(nullptr) {}
};

class TreeNode {
public:
  int val;
  TreeNode *left, *right;
  TreeNode(int v) : val(v), left(nullptr), right(nullptr) {}
};
\end{lstlisting}

\algotitle{1.8 Contest Complexity \& Constants Cheat Sheet}
\algodesc{Rules of thumb for 1.0 second time limits ($\approx 10^8$ operations).}
\begin{itemize}[leftmargin=*,noitemsep,topsep=1pt]
  \item $N \le 10$: $O(N!) \approx 3.6 \times 10^6$
  \item $N \le 20$: $O(2^N) \approx 10^6$
  \item $N \le 500$: $O(N^3) \approx 1.25 \times 10^8$
  \item $N \le 5000$: $O(N^2) \approx 2.5 \times 10^7$
  \item $N \le 10^6$: $O(N \log N) \approx 2 \times 10^7$
  \item $N \le 10^8$: $O(N)$ linear
  \item $2^{10} \approx 10^3$, $2^{20} \approx 10^6$, $2^{30} \approx 10^9$, $2^{60} \approx 10^{18}$
\end{itemize}



\catsection{2. Mathematics \& Basic Algebra}

\algotitle{2.1 Series \& Sums Formulas}
\algodesc{Closed form expressions for standard summations and power series.}
\begin{itemize}[leftmargin=*,noitemsep,topsep=1pt]
  \item Arithmetic Progression: $S_n = \frac{n}{2}(2a + (n-1)d) = \frac{n}{2}(a_1 + a_n)$
  \item Geometric Progression: $\sum_{i=0}^n c^i = \frac{c^{n+1}-1}{c-1}$, $\sum_{i=a}^b c^i = \frac{c^{b+1}-c^a}{c-1} \ (c \ne 1)$
  \item Infinite GP: $\sum_{i=0}^\infty c^i = \frac{1}{1-c} \ (|c| < 1)$
  \item Sum of first $n$ integers: $\sum_{i=1}^n i = \frac{n(n+1)}{2}$
  \item Sum of squares: $\sum_{i=1}^n i^2 = \frac{n(n+1)(2n+1)}{6}$
  \item Sum of cubes: $\sum_{i=1}^n i^3 = \left[\frac{n(n+1)}{2}\right]^2$
  \item Sum of 4th powers: $\sum_{i=1}^n i^4 = \frac{n(n+1)(2n+1)(3n^2+3n-1)}{30}$
  \item Harmonic Numbers: $H_n = \sum_{i=1}^n \frac{1}{i} \approx \ln n + \gamma \ (\gamma \approx 0.57721566)$
  \item Weighted Harmonic: $\sum_{i=1}^n i H_i = \frac{n(n+1)}{2}H_n - \frac{n(n-1)}{4}$
\end{itemize}

\algotitle{2.2 Algebraic Identities}
\algodesc{Factorizations and expansions frequently tested in algebraic manipulations.}
\begin{itemize}[leftmargin=*,noitemsep,topsep=1pt]
  \item $(a \pm b)^2 = a^2 \pm 2ab + b^2 \quad \vert \quad a^2 - b^2 = (a-b)(a+b)$
  \item $(a \pm b)^3 = a^3 \pm 3a^2b + 3ab^2 \pm b^3 \quad \vert \quad a^3 \pm b^3 = (a \pm b)(a^2 \mp ab + b^2)$
  \item $(a + b + c)^2 = a^2 + b^2 + c^2 + 2(ab + bc + ca)$
  \item $a^4 + 4b^4 = (a^2 - 2ab + 2b^2)(a^2 + 2ab + 2b^2)$ \ (Sophie Germain)
  \item $a^n - b^n = (a-b)(a^{n-1} + a^{n-2}b + \dots + b^{n-1})$
  \item $a^n + b^n = (a+b)(a^{n-1} - a^{n-2}b + \dots + b^{n-1})$ \ (for odd $n$)
\end{itemize}

\algotitle{2.3 Classical Inequalities}
\algodesc{Inequalities used for bounds, binary searches, and proofs.}
\begin{itemize}[leftmargin=*,noitemsep,topsep=1pt]
  \item AM-GM: $\frac{x_1 + x_2 + \dots + x_n}{n} \ge \sqrt[n]{x_1 x_2 \dots x_n}$ \ (equality iff all $x_i$ equal)
  \item Cauchy-Schwarz: $\left(\sum a_i^2\right)\left(\sum b_i^2\right) \ge \left(\sum a_i b_i\right)^2$
  \item Triangle: $\big| |a| - |b| \big| \le |a \pm b| \le |a| + |b|$
  \item Jensen's: $f\left(\frac{\sum x_i}{n}\right) \le \frac{\sum f(x_i)}{n}$ for convex $f$ ($f'' \ge 0$)
  \item Bernoulli: $(1+x)^n \ge 1 + nx$ for $n \ge 0, x > -1$
  \item Chebyshev: $\frac{1}{n}\sum a_i b_i \ge \left(\frac{1}{n}\sum a_i\right)\left(\frac{1}{n}\sum b_i\right)$ if sorted similarly
\end{itemize}

\algotitle{2.4 Probability \& Expected Value}
\algodesc{Rules of probability, linearity of expectation, and statistical bounds.}
\begin{itemize}[leftmargin=*,noitemsep,topsep=1pt]
  \item $P(A \cup B) = P(A) + P(B) - P(A \cap B)$
  \item Bayes' Theorem: $P(A|B) = \frac{P(B|A)P(A)}{P(B)}$
  \item Linearity of Expectation: $E[X + Y] = E[X] + E[Y]$ (ALWAYS holds, even if $X, Y$ are dependent!)
  \item Independent Variables: $E[X \cdot Y] = E[X] \cdot E[Y]$
  \item Variance: $\operatorname{Var}(X) = E[X^2] - (E[X])^2$, $\operatorname{Var}(aX+b) = a^2\operatorname{Var}(X)$
  \item Markov: $P(X \ge a) \le \frac{E[X]}{a}$ \ ($X \ge 0, a > 0$)
  \item Chebyshev: $P(|X - \mu| \ge k\sigma) \le \frac{1}{k^2}$
\end{itemize}

\algotitle{2.5 Binary Exponentiation}
\algodesc{Computes $a^b \pmod M$ in $O(\log b)$ time.}
\begin{lstlisting}
int binexp(int a, int b, int m = M) {
  int res = 1 % m; a %= m;
  while (b > 0) {
    if (b & 1) res = (__int128)res * a % m;
    a = (__int128)a * a % m;
    b >>= 1;
  }
  return res;
}
\end{lstlisting}

\algotitle{2.6 Matrix Exponentiation}
\algodesc{Solves linear recurrences and counts graph paths of length $K$ in $O(N^3 \log K)$.}
\begin{lstlisting}
template<size_t N>
struct Matrix {
  int mat[N][N];
  Matrix(int v = 0) {
    for (size_t i = 0; i < N; i++)
      for (size_t j = 0; j < N; j++)
        mat[i][j] = (i == j ? v : 0);
  }
  Matrix operator*(const Matrix& other) const {
    Matrix res(0);
    for (size_t i = 0; i < N; i++) {
      for (size_t k = 0; k < N; k++) {
        if (!mat[i][k]) continue;
        for (size_t j = 0; j < N; j++) {
          res.mat[i][j] = (res.mat[i][j] + (__int128)mat[i][k] * other.mat[k][j]) % M;
        }
      }
    }
    return res;
  }
};
template<size_t N>
Matrix<N> mat_expo(Matrix<N> a, int p) {
  Matrix<N> res(1);
  while (p > 0) {
    if (p & 1) res = res * a;
    a = a * a;
    p >>= 1;
  }
  return res;
}
// Fibonacci: F[n] = F[n-1] + F[n-2], F[0]=0, F[1]=1
int fib_mat(int n) {
  if (n == 0) return 0;
  Matrix<2> T;
  T.mat[0][0] = 1; T.mat[0][1] = 1;
  T.mat[1][0] = 1; T.mat[1][1] = 0;
  T = mat_expo(T, n - 1);
  return T.mat[0][0]; // F[n]
}
\end{lstlisting}

\algotitle{2.7 Basis Vector (XOR Basis)}
\algodesc{Maintains linear basis of vectors over $\mathbb{F}_2$. Supports max XOR subset, $k$-th XOR.}
\begin{lstlisting}
template<typename T = int, int LG = 60>
struct Basis {
  T basis[LG];
  int sz;
  Basis() { reset(); }
  void reset() { fill(basis, basis + LG, 0); sz = 0; }
  bool insert(T mask) {
    for (int i = LG - 1; i >= 0; --i) {
      if ((mask >> i) & 1) {
        if (basis[i]) mask ^= basis[i];
        else { basis[i] = mask; ++sz; return true; }
      }
    }
    return false;
  }
  T maxXor() {
    T ans = 0;
    for (int i = LG - 1; i >= 0; --i) ans = max(ans, ans ^ basis[i]);
    return ans;
  }
  bool can(T x) {
    for (int i = LG - 1; i >= 0; --i) {
      if ((x >> i) & 1) {
        if (!basis[i]) return false;
        x ^= basis[i];
      }
    }
    return true;
  }
  T kth(T k) { // 1-indexed k-th smallest XOR
    if (k < 1 || k > (1ULL << sz)) return -1;
    T mask = 0, tot = 1ULL << sz;
    for (int i = LG - 1; i >= 0; --i) {
      if (basis[i]) {
        T low = tot / 2;
        if ((low < k && !((mask >> i) & 1)) || (low >= k && ((mask >> i) & 1)))
          mask ^= basis[i];
        if (low < k) k -= low;
        tot /= 2;
      }
    }
    return mask;
  }
};
\end{lstlisting}

\algotitle{2.8 Linear Combination Checking (Bézout's)}
\algodesc{Check if integer $N$ can be formed by $a \cdot x + b \cdot y + c \cdot z + \dots$}
\begin{itemize}[leftmargin=*,noitemsep,topsep=1pt]
  \item $N$ can be represented iff $N \bmod \gcd(a_1, a_2, \dots, a_k) == 0$.
  \item Frobenius Coin Problem: For 2 coprime integers $a, b$, the largest number that CANNOT be formed is $a \cdot b - a - b$. Exactly $\frac{(a-1)(b-1)}{2}$ numbers cannot be formed.
\end{itemize}



\catsection{3. Number Theory \& Divisors}

\algotitle{3.1 GCD \& LCM Fundamental Properties}
\algodesc{Mathematical identities connecting GCD, LCM, divisibility, and lattice lines.}
\begin{itemize}[leftmargin=*,noitemsep,topsep=1pt]
  \item $\gcd(a, b) \times \operatorname{lcm}(a, b) = |a \cdot b|$
  \item $\gcd(a, b) = \gcd(a, b \bmod a) = \gcd(a, b - a)$
  \item $\gcd(k \cdot a, k \cdot b) = k \cdot \gcd(a, b)$
  \item $\gcd(a, b) = 1 \iff \gcd(a + b, a \cdot b) = 1$
  \item If $a \mid (b \cdot c)$ and $\gcd(a, b) = d$, then $\frac{a}{d} \mid c$.
  \item Exponents: $\gcd(a^m - 1, a^n - 1) = a^{\gcd(m, n)} - 1$
  \item Fibonacci GCD: $\gcd(F_m, F_n) = F_{\gcd(m, n)}$
  \item Lattice segment: The number of strict lattice points between $(0, 0)$ and $(a, b)$ is $\gcd(|a|, |b|) - 1$.
\end{itemize}

\algotitle{3.2 Extended Euclidean Algorithm \& Diophantine}
\algodesc{Solves $a \cdot x + b \cdot y = \gcd(a, b)$ and linear Diophantine $a \cdot x + b \cdot y = c$.}
\begin{lstlisting}
int extgcd(int a, int b, int &x, int &y) {
  if (b == 0) { x = 1; y = 0; return a; }
  int x1, y1;
  int d = extgcd(b, a % b, x1, y1);
  x = y1;
  y = x1 - y1 * (a / b);
  return d;
}
bool find_any_solution(int a, int b, int c, int &x0, int &y0, int &g) {
  g = extgcd(abs(a), abs(b), x0, y0);
  if (c % g != 0) return false;
  x0 *= c / g; y0 *= c / g;
  if (a < 0) x0 = -x0;
  if (b < 0) y0 = -y0;
  // All solutions: x = x0 + k*(b/g), y = y0 - k*(a/g)
  return true;
}
\end{lstlisting}

\algotitle{3.3 Modular Inverse}
\algodesc{Fermat's theorem for prime mod $M$; Extended Euclidean for arbitrary coprime $m$.}
\begin{lstlisting}
// Prime mod: a^(m-2) mod m
int modInvPrime(int a, int m = M) { return binexp(a, m - 2, m); }

// Arbitrary mod (requires gcd(a, m) == 1):
int modInv(int a, int m) {
  int x, y;
  int g = extgcd(a, m, x, y);
  if (g != 1) return -1; // No inverse exists
  return (x % m + m) % m;
}
\end{lstlisting}

\algotitle{3.4 Chinese Remainder Theorem (General CRT)}
\algodesc{Merges $x \equiv a_1 \pmod{m_1}$ and $x \equiv a_2 \pmod{m_2}$. Supports non-coprime moduli.}
\begin{lstlisting}
pair<int, int> CRT(int a1, int m1, int a2, int m2) {
  int x, y;
  int g = extgcd(m1, m2, x, y);
  if ((a2 - a1) % g != 0) return {-1, -1}; // No solution
  int lcm = (m1 / g) * m2;
  int t = ((__int128)(a2 - a1) / g * x) % (m2 / g);
  int ans = (a1 + (__int128)m1 * t) % lcm;
  if (ans < 0) ans += lcm;
  return {ans, lcm}; // x === ans (mod lcm)
}
\end{lstlisting}

\algotitle{3.5 Sieve of Eratosthenes}
\algodesc{Finds all primes up to $N$ in $O(N \log \log N)$.}
\begin{lstlisting}
vector<int> primes;
vector<bool> is_prime(N + 1, true);
void sieve() {
  is_prime[0] = is_prime[1] = false;
  for (int i = 2; i * i <= N; i++) {
    if (is_prime[i]) {
      for (int j = i * i; j <= N; j += i) is_prime[j] = false;
    }
  }
  for (int i = 2; i <= N; i++) if (is_prime[i]) primes.push_back(i);
}
\end{lstlisting}

\algotitle{3.6 Smallest Prime Factor (SPF) \& Linear Sieve}
\algodesc{Computes SPF and primes in $O(N)$ time. Enables $O(\log N)$ factorization.}
\begin{lstlisting}
vector<int> spf(N + 1), primes;
void linear_sieve() {
  for (int i = 2; i <= N; i++) {
    if (!spf[i]) { spf[i] = i; primes.push_back(i); }
    for (int p : primes) {
      if (p > spf[i] || i * p > N) break;
      spf[i * p] = p;
    }
  }
}
// Factorization via SPF in O(log N):
vector<pair<int, int>> factorize_spf(int n) {
  vector<pair<int, int>> factors;
  while (n > 1) {
    int p = spf[n], c = 0;
    while (n % p == 0) { n /= p; c++; }
    factors.push_back({p, c});
  }
  return factors;
}
\end{lstlisting}

\algotitle{3.7 Prime Factorization in $O(\sqrt{N})$}
\algodesc{Standard factorization for isolated large numbers up to $10^{14}$.}
\begin{lstlisting}
vector<pair<int, int>> factorize(int n) {
  vector<pair<int, int>> factors;
  for (int d = 2; d * d <= n; d++) {
    if (n % d == 0) {
      int count = 0;
      while (n % d == 0) { n /= d; count++; }
      factors.push_back({d, count});
    }
  }
  if (n > 1) factors.push_back({n, 1});
  return factors;
}
\end{lstlisting}

\algotitle{3.8 Segmented Sieve}
\algodesc{Finds all primes in range $[L, R]$ where $R \le 10^{12}, R - L \le 10^7$.}
\begin{lstlisting}
vector<int> segmented_sieve(int L, int R) {
  int lim = sqrt(R);
  vector<bool> mark(lim + 1, false);
  vector<int> small_primes;
  for (int i = 2; i <= lim; i++) {
    if (!mark[i]) {
      small_primes.push_back(i);
      for (int j = i * i; j <= lim; j += i) mark[j] = true;
    }
  }
  vector<bool> isPrime(R - L + 1, true);
  for (int p : small_primes) {
    int sm = (L + p - 1) / p * p;
    for (int j = max(p * p, sm); j <= R; j += p) isPrime[j - L] = false;
  }
  if (L <= 1) for (int i = L; i <= min(1LL, R); i++) isPrime[i - L] = false;
  vector<int> res;
  for (int i = 0; i <= R - L; i++) if (isPrime[i]) res.push_back(L + i);
  return res;
}
\end{lstlisting}

\algotitle{3.9 Divisor Formulas (Complete Theory \& Bounds)}
\algodesc{All formulas for prime factorization $n = p_1^{a_1} p_2^{a_2} \dots p_k^{a_k}$.}
\begin{itemize}[leftmargin=*,noitemsep,topsep=1pt]
  \item Number of Divisors: $\tau(n) = d(n) = \prod_{i=1}^k (a_i + 1)$
  \item Sum of Divisors: $\sigma(n) = \prod_{i=1}^k \frac{p_i^{a_i+1} - 1}{p_i - 1}$
  \item Product of Divisors: $\prod_{d \mid n} d = n^{\tau(n) / 2}$
  \item Odd / Even Divisors: If $n = 2^a \cdot m$ ($m$ odd), then number of odd divisors is $\tau(m)$, and number of even divisors is $a \cdot \tau(m)$.
  \item Sum of Powers of Divisors: $\sigma_k(n) = \prod_{i=1}^k \frac{p_i^{(a_i+1)k} - 1}{p_i^k - 1}$
  \item Dirichlet Hyperbola Method: $\sum_{i=1}^N d(i) = \sum_{i=1}^N \lfloor N/i \rfloor = 2 \sum_{i=1}^{\lfloor\sqrt{N}\rfloor} \lfloor N/i \rfloor - \lfloor\sqrt{N}\rfloor^2$ \ ($O(\sqrt{N})$)
  \item Max Number of Divisors Bounds (Highly Composite Numbers):
    \begin{itemize}[noitemsep]
      \item $N \le 10^3: \max d(n) = 32 \ (n = 840)$
      \item $N \le 10^6: \max d(n) = 240 \ (n = 720720)$
      \item $N \le 10^9: \max d(n) = 1344 \ (n = 735134400)$
      \item $N \le 10^{12}: \max d(n) = 6720 \ (n = 963761198400)$
      \item $N \le 10^{18}: \max d(n) = 103680 \dots 4.35 \times 10^7$
    \end{itemize}
\end{itemize}

\algotitle{3.10 Sum of Floor($N/i$) in $O(\sqrt{N})$}
\algodesc{Groups equal values of $\lfloor N/i \rfloor$ into blocks $[l, r]$.}
\begin{lstlisting}
int sum_floor(int n) {
  int ans = 0;
  for (int l = 1, r; l <= n; l = r + 1) {
    r = n / (n / l);
    ans += (r - l + 1) * (n / l);
  }
  return ans;
}
\end{lstlisting}

\algotitle{3.11 Euler's Totient Function (ETF / $\phi(n)$)}
\algodesc{Counts integers in $[1, n]$ coprime to $n$. $\phi(n) = n \prod_{p \mid n} (1 - 1/p)$.}
\begin{lstlisting}
// O(sqrt(N)) for single query:
int phi(int n) {
  int res = n;
  for (int i = 2; i * i <= n; i++) {
    if (n % i == 0) {
      while (n % i == 0) n /= i;
      res -= res / i;
    }
  }
  if (n > 1) res -= res / n;
  return res;
}
// Sieve for all phi[1..N] in O(N):
vector<int> phi_arr(N + 1);
void sieve_phi() {
  iota(phi_arr.begin(), phi_arr.end(), 0);
  for (int i = 2; i <= N; i++) {
    if (phi_arr[i] == i) {
      for (int j = i; j <= N; j += i) phi_arr[j] -= phi_arr[j] / i;
    }
  }
}
\end{lstlisting}
\algodesc{ETF Key Properties: $\sum_{d \mid n} \phi(d) = n$; Sum of coprimes $\sum_{1 \le i \le n, \gcd(i,n)=1} i = \frac{n \cdot \phi(n)}{2}$. Euler's Totient Theorem: $a^{\phi(m)} \equiv 1 \pmod m$ for $\gcd(a, m) = 1$. General Power Tower Reduction: $a^b \equiv a^{(b \bmod \phi(m)) + \phi(m)} \pmod m$ for $b \ge \phi(m)$.}

\algotitle{3.12 Miller-Rabin Primality Test (Deterministic 64-bit)}
\algodesc{Checks if $N < 2^{64}$ is prime in $O(\log^3 N)$ using deterministic witness bases.}
\begin{lstlisting}
using u64 = uint64_t;
using u128 = __uint128_t;
u64 bin_pow(u64 a, u64 d, u64 mod) {
  u64 res = 1; a %= mod;
  while (d > 0) {
    if (d & 1) res = (u128)res * a % mod;
    a = (u128)a * a % mod;
    d >>= 1;
  }
  return res;
}
bool miller_rabin(u64 n) {
  if (n < 2) return false;
  if (n == 2 || n == 3) return true;
  if (n % 2 == 0) return false;
  u64 d = n - 1; int s = 0;
  while ((d & 1) == 0) { d >>= 1; s++; }
  static const u64 bases[] = {2, 325, 9375, 28178, 450775, 9780504, 1795265022};
  for (u64 a : bases) {
    if (a % n == 0) continue;
    u64 x = bin_pow(a, d, n);
    if (x == 1 || x == n - 1) continue;
    bool composite = true;
    for (int r = 1; r < s; r++) {
      x = (u128)x * x % n;
      if (x == n - 1) { composite = false; break; }
    }
    if (composite) return false;
  }
  return true;
}
\end{lstlisting}

\algotitle{3.13 Pollard's Rho (Fast 64-bit Factorization)}
\algodesc{Finds a non-trivial factor of $N$ in expected $O(N^{1/4})$ time.}
\begin{lstlisting}
int pollard_rho(int n) {
  if (n % 2 == 0) return 2;
  int x = 2, y = 2, d = 1, c = 1;
  auto f = [&](int val) { return ((__int128)val * val + c) % n; };
  while (d == 1) {
    x = f(x); y = f(f(y));
    d = std::gcd(abs(x - y), n);
    if (d == n) { x = rand() % (n - 2) + 2; y = x; c = rand() % (n - 1) + 1; d = 1; }
  }
  return d;
}
\end{lstlisting}

\algotitle{3.14 Floor Sum (AtCoder ACL)}
\algodesc{Computes $\sum_{i=0}^{n-1} \lfloor \frac{a \cdot i + b}{m} \rfloor$ in $O(\log m)$ via Euclidean division.}
\begin{lstlisting}
int floor_sum(int n, int m, int a, int b) {
  int ans = 0;
  if (a >= m) { ans += (n - 1) * n * (a / m) / 2; a %= m; }
  if (b >= m) { ans += n * (b / m); b %= m; }
  int y_max = (a * n + b) / m;
  if (y_max == 0) return ans;
  ans += floor_sum(y_max, a, m, (a * n + b) % m);
  return ans;
}
\end{lstlisting}

\algotitle{3.15 Discrete Logarithm (Baby-Step Giant-Step)}
\algodesc{Finds smallest $x \ge 0$ such that $a^x \equiv b \pmod m$ with $\gcd(a, m) = 1$ in $O(\sqrt{m})$.}
\begin{lstlisting}
int baby_step_giant_step(int a, int b, int m) {
  a %= m; b %= m;
  if (b == 1 || m == 1) return 0;
  int n = sqrt(m) + 1;
  unordered_map<int, int> tbl;
  int cur = b;
  for (int q = 0; q < n; q++) {
    tbl[cur] = q;
    cur = (__int128)cur * a % m;
  }
  int an = binexp(a, n, m);
  cur = 1;
  for (int p = 1; p <= n + 1; p++) {
    cur = (__int128)cur * an % m;
    if (tbl.count(cur)) {
      int ans = p * n - tbl[cur];
      return ans;
    }
  }
  return -1;
}
\end{lstlisting}

\algotitle{3.16 Möbius Inversion Formula}
\algodesc{$\mu(n) = (-1)^k$ if $n$ is square-free with $k$ prime factors, $0$ if $p^2 \mid n$.}
\begin{itemize}[leftmargin=*,noitemsep,topsep=1pt]
  \item $\sum_{d \mid n} \mu(d) = [n == 1]$
  \item Inversion: $g(n) = \sum_{d \mid n} f(d) \iff f(n) = \sum_{d \mid n} \mu(d) g(n / d)$
  \item Co-primality count: $[ \gcd(i, j) == 1 ] = \sum_{d \mid \gcd(i, j)} \mu(d)$
  \item GCD pairs: $\sum_{1 \le i, j \le N} [\gcd(i, j) = 1] = \sum_{d=1}^N \mu(d) \lfloor N / d \rfloor^2$
\end{itemize}

\algotitle{3.17 Lifting The Exponent (LTE)}
\algodesc{$\nu_p(x)$ is exponent of highest power of prime $p$ dividing $x$.}
\begin{itemize}[leftmargin=*,noitemsep,topsep=1pt]
  \item Odd Prime ($p \ge 3$), $p \mid (a - b)$, $p \nmid a, p \nmid b$: $\nu_p(a^n - b^n) = \nu_p(a - b) + \nu_p(n)$
  \item Odd Prime ($p \ge 3$), $n$ odd, $p \mid (a + b)$: $\nu_p(a^n + b^n) = \nu_p(a + b) + \nu_p(n)$
  \item Even Prime ($p = 2$), $a, b$ odd:
    $\nu_2(a^n - b^n) = \nu_2(a - b) + \nu_2(a + b) + \nu_2(n) - 1$ \ (for even $n$)
\end{itemize}

\algotitle{3.18 Fibonacci Properties \& Fast Doubling}
\algodesc{Computes $\{F_n, F_{n+1}\}$ in $O(\log n)$ using $F_{2k} = F_k(2F_{k+1} - F_k), F_{2k+1} = F_k^2 + F_{k+1}^2$.}
\begin{lstlisting}
pair<int, int> fib_fast(int n) {
  if (n == 0) return {0, 1};
  auto [a, b] = fib_fast(n >> 1);
  int c = (a * (2 * b - a + M) % M) % M;
  int d = (a * a + b * b) % M;
  if (n & 1) return {d, (c + d) % M};
  return {c, d};
}
\end{lstlisting}



\catsection{4. Combinatorics \& Discrete Math}

\algotitle{4.1 Permutations, Combinations \& Modular nCr}
\algodesc{Precomputes factorials and inverse factorials up to $N$ in $O(N)$ for $O(1)$ queries.}
\begin{lstlisting}
vector<int> fact(N + 5), ifact(N + 5);
void precompute_nCr() {
  fact[0] = 1;
  for (int i = 1; i <= N; i++) fact[i] = fact[i - 1] * i % M;
  ifact[N] = binexp(fact[N], M - 2, M);
  for (int i = N - 1; i >= 0; i--) ifact[i] = ifact[i + 1] * (i + 1) % M;
}
int nCr(int n, int r) {
  if (r < 0 || r > n) return 0;
  return fact[n] * ifact[r] % M * ifact[n - r] % M;
}
int nPr(int n, int r) {
  if (r < 0 || r > n) return 0;
  return fact[n] * ifact[n - r] % M;
}
\end{lstlisting}

\algotitle{4.2 Binomial Coefficient Properties}
\algodesc{Core identities for manipulating combinations in CP math problems.}
\begin{itemize}[leftmargin=*,noitemsep,topsep=1pt]
  \item Symmetry: $\binom{n}{k} = \binom{n}{n-k}$
  \item Pascal's Rule: $\binom{n}{k} = \binom{n-1}{k-1} + \binom{n-1}{k}$
  \item Absorption (Committee): $k \binom{n}{k} = n \binom{n-1}{k-1}$
  \item Total Subsets: $\sum_{k=0}^n \binom{n}{k} = 2^n$
  \item Alternating Sum: $\sum_{k=0}^n (-1)^k \binom{n}{k} = 0$
  \item Even/Odd Sums: $\sum_{k \text{ even}} \binom{n}{k} = \sum_{k \text{ odd}} \binom{n}{k} = 2^{n-1}$
  \item Vandermonde's Convolution: $\sum_{k=0}^r \binom{m}{k}\binom{n}{r-k} = \binom{m+n}{r}$
  \item Sum of Squares: $\sum_{k=0}^n \binom{n}{k}^2 = \binom{2n}{n}$
  \item Weighted Sum: $\sum_{k=1}^n k \binom{n}{k} = n \cdot 2^{n-1}$
\end{itemize}

\algotitle{4.3 Hockey-Stick Identity (Pyramid Formula for nCr)}
\algodesc{Summing along a diagonal or column in Pascal's triangle.}
\begin{itemize}[leftmargin=*,noitemsep,topsep=1pt]
  \item Column Sum: $\sum_{i=r}^n \binom{i}{r} = \binom{n+1}{r+1}$
  \item Diagonal Sum: $\sum_{i=0}^k \binom{n+i}{i} = \binom{n+k+1}{k}$
  \item Pascal's Pyramid (3D / Multinomial):
    $\binom{n}{i, j, k} = \frac{n!}{i! j! k!} \quad (i + j + k = n)$
    Each entry is the sum of the 3 adjacent entries in the layer above:
    $\binom{n}{i, j, k} = \binom{n-1}{i-1, j, k} + \binom{n-1}{i, j-1, k} + \binom{n-1}{i, j, k-1}$
\end{itemize}

\algotitle{4.4 Stars and Bars (Ballot Problems)}
\algodesc{Distributing $n$ identical objects into $k$ distinct containers.}
\begin{itemize}[leftmargin=*,noitemsep,topsep=1pt]
  \item Non-negative integers ($x_i \ge 0$): $\binom{n + k - 1}{k - 1}$
  \item Strictly positive integers ($x_i \ge 1$): $\binom{n - 1}{k - 1}$
  \item Lower bounded ($x_i \ge c_i$): Let $n' = n - \sum c_i$; ways $= \binom{n' + k - 1}{k - 1}$
\end{itemize}

\algotitle{4.5 Catalan Numbers (Formulas \& CP Explanations)}
\algodesc{Sequence: $1, 1, 2, 5, 14, 42, 132, 429, 1430, 4862, \dots$}
\begin{itemize}[leftmargin=*,noitemsep,topsep=1pt]
  \item Explicit Formula: $C_n = \frac{1}{n+1}\binom{2n}{n} = \binom{2n}{n} - \binom{2n}{n+1} = \frac{(2n)!}{(n+1)! n!}$
  \item Recurrence: $C_0 = 1, \quad C_{n+1} = \sum_{i=0}^n C_i C_{n-i} = \frac{2(2n+1)}{n+2} C_n$
  \item CP Interpretation \& Applications:
    \begin{enumerate}[noitemsep]
      \item Balanced Parentheses: Number of valid expressions with $n$ pairs of `(' and `)'.
      \item Dyck Paths: Monotonic grid walks from $(0, 0)$ to $(n, n)$ that never cross strictly above the diagonal $y = x$.
      \item Polygon Triangulation: Ways to divide a convex polygon of $n+2$ vertices into $n$ triangles with non-intersecting diagonals.
      \item Full Binary Trees: Number of distinct rooted binary trees with $n+1$ leaves ($n$ internal nodes).
      \item Non-intersecting Chords: Ways to pair $2n$ points on a circle with $n$ non-intersecting chords.
      \item Stack Permutations: Permutations of length $n$ that can be sorted using a single stack (avoiding pattern 231).
    \end{enumerate}
\end{itemize}
\begin{lstlisting}
int catalan(int n) {
  return nCr(2 * n, n) * modInvPrime(n + 1) % M;
}
\end{lstlisting}

\algotitle{4.6 Derangements $D(n)$}
\algodesc{Permutations where no element appears in its original position ($\pi(i) \ne i$).}
\begin{itemize}[leftmargin=*,noitemsep,topsep=1pt]
  \item Sequence: $0, 1, 2, 9, 44, 265, 1854, 14833, \dots$
  \item Recurrence: $D(n) = (n - 1)(D(n - 1) + D(n - 2)) = n D(n - 1) + (-1)^n$
  \item Closed form: $D(n) = n! \sum_{k=0}^n \frac{(-1)^k}{k!} \approx \left[\frac{n!}{e}\right]$
\end{itemize}

\algotitle{4.7 Stirling Numbers of 1st and 2nd Kind}
\algodesc{Cycle arrangements and set partitions.}
\begin{itemize}[leftmargin=*,noitemsep,topsep=1pt]
  \item 1st Kind $\left[\begin{matrix}n\\k\end{matrix}\right]$: Permutations of $n$ elements with exactly $k$ disjoint cycles.
    $\left[\begin{matrix}n\\k\end{matrix}\right] = \left[\begin{matrix}n-1\\k-1\end{matrix}\right] + (n-1)\left[\begin{matrix}n-1\\k\end{matrix}\right], \quad \sum_{k=0}^n \left[\begin{matrix}n\\k\end{matrix}\right] = n!$
  \item 2nd Kind $\left\{\begin{matrix}n\\k\end{matrix}\right\}$: Ways to partition $n$ distinct elements into $k$ non-empty sets.
    $\left\{\begin{matrix}n\\k\end{matrix}\right\} = \left\{\begin{matrix}n-1\\k-1\end{matrix}\right\} + k\left\{\begin{matrix}n-1\\k\end{matrix}\right\}, \quad \left\{\begin{matrix}n\\k\end{matrix}\right\} = \frac{1}{k!}\sum_{i=0}^k (-1)^{k-i}\binom{k}{i}i^n$
  \item Bell Numbers $B(n) = \sum_{k=0}^n \left\{\begin{matrix}n\\k\end{matrix}\right\}$: Total ways to partition $n$ distinct elements.
\end{itemize}

\algotitle{4.8 Lucas' Theorem}
\algodesc{Computes $\binom{n}{r} \pmod p$ for large $n, r$ when modulo $p$ is a small prime.}
\begin{lstlisting}
int lucas(int n, int r, int p) {
  if (r == 0) return 1;
  int ni = n % p, ri = r % p;
  if (ni < ri) return 0;
  return lucas(n / p, r / p, p) * nCr(ni, ri) % p;
}
\end{lstlisting}

\algotitle{4.9 Inclusion-Exclusion Principle}
\algodesc{Computes size of union of properties $|\bigcup A_i| = \sum_{k=1}^n (-1)^{k-1} \sum_{|J|=k} |\bigcap_{j \in J} A_j|$.}
\begin{lstlisting}
// Count elements satisfying none of n conditions
template<typename Func>
int inclusion_exclusion(int n, Func f) {
  int ans = 0;
  for (int mask = 0; mask < (1 << n); mask++) {
    int sign = (__builtin_popcount(mask) & 1) ? -1 : 1;
    ans += sign * f(mask);
  }
  return ans;
}
\end{lstlisting}

\algotitle{4.10 Burnside's Lemma (Colorings Under Symmetry)}
\algodesc{Number of orbits under group actions: $|X/G| = \frac{1}{|G|} \sum_{g \in G} |X^g|$ where $|X^g| = k^{c(g)}$.}
\begin{itemize}[leftmargin=*,noitemsep,topsep=1pt]
  \item Necklaces with $n$ beads and $k$ colors under rotation: $\frac{1}{n} \sum_{i=0}^{n-1} k^{\gcd(i, n)}$
  \item Bracelets (rotation + reflection): Add reflection symmetries to group size.
\end{itemize}



\catsection{5. Data Structures}

\algotitle{5.1 Policy Based Data Structures (PBDS Ordered Set)}
\algodesc{Supports find\_by\_order(k) and order\_of\_key(x) in $O(\log N)$. Multiset uses pair<T, int> to avoid UB.}
\begin{lstlisting}
#include <ext/pb_ds/assoc_container.hpp>
#include <ext/pb_ds/tree_policy.hpp>
using namespace __gnu_pbds;

template <class T, class Cmp = less<T>>
using ordered_set = tree<T, null_type, Cmp, rb_tree_tag, tree_order_statistics_node_update>;

// Safe Multiset without UB:
ordered_set<pair<int, int>> ost;
int timer = 0;
void insert_val(int x) { ost.insert({x, ++timer}); }
void erase_val(int x) {
  auto it = ost.lower_bound({x, 0});
  if (it != ost.end() && it->first == x) ost.erase(it);
}
int order_of(int x) { return ost.order_of_key({x, 0}); } // count < x
int kth_element(int k) { return ost.find_by_order(k)->first; } // 0-indexed
\end{lstlisting}

\algotitle{5.2 Disjoint Set Union (DSU)}
\algodesc{Standard DSU with path compression and union by size. Tracks component count and sizes.}
\begin{lstlisting}
struct DSU {
  vector<int> p, sz;
  int comps;
  DSU(int n) : p(n + 1), sz(n + 1, 1), comps(n) {
    iota(p.begin(), p.end(), 0);
  }
  int find(int x) {
    return p[x] == x ? x : p[x] = find(p[x]);
  }
  bool unite(int a, int b) {
    a = find(a); b = find(b);
    if (a == b) return false;
    if (sz[a] < sz[b]) swap(a, b);
    p[b] = a;
    sz[a] += sz[b];
    comps--;
    return true;
  }
  int size(int x) { return sz[find(x)]; }
};
\end{lstlisting}

\algotitle{5.3 Rollback DSU (Undoable Disjoint Sets)}
\algodesc{No path compression; union by rank only. Used for offline dynamic connectivity and divide-and-conquer.}
\begin{lstlisting}
struct RollbackDSU {
  vector<int> p, sz;
  vector<pair<int, int>> history;
  RollbackDSU(int n) : p(n + 1), sz(n + 1, 1) {
    iota(p.begin(), p.end(), 0);
  }
  int find(int x) {
    while (x != p[x]) x = p[x];
    return x;
  }
  bool unite(int a, int b) {
    a = find(a); b = find(b);
    if (a == b) return false;
    if (sz[a] < sz[b]) swap(a, b);
    history.push_back({b, sz[a]});
    p[b] = a;
    sz[a] += sz[b];
    return true;
  }
  int time() { return history.size(); }
  void rollback(int t) {
    while ((int)history.size() > t) {
      auto [b, old_sz_a] = history.back(); history.pop_back();
      int a = p[b];
      sz[a] = old_sz_a;
      p[b] = b;
    }
  }
};
\end{lstlisting}

\algotitle{5.4 Fenwick Tree (Binary Indexed Tree)}
\algodesc{1-indexed point update, range query, and $O(\log N)$ binary lifting search for prefix sums.}
\begin{lstlisting}
struct Fenwick {
  int n;
  vector<int> tree;
  Fenwick(int n) : n(n), tree(n + 1, 0) {}
  void add(int i, int delta) {
    for (; i <= n; i += i & -i) tree[i] += delta;
  }
  int query(int i) {
    int sum = 0;
    for (; i > 0; i -= i & -i) sum += tree[i];
    return sum;
  }
  int query(int l, int r) { return query(r) - query(l - 1); }
  int find_kth(int k) { // Find smallest index with prefix sum >= k
    int idx = 0;
    for (int j = 1 << 20; j > 0; j >>= 1) {
      if (idx + j <= n && tree[idx + j] < k) {
        idx += j;
        k -= tree[idx];
      }
    }
    return idx + 1;
  }
};
\end{lstlisting}

\algotitle{5.5 Range Update \& Range Query Fenwick (2-BIT Trick)}
\algodesc{Supports range add and range sum using $p(i) = i \cdot \text{bit}_1(i) - \text{bit}_2(i)$.}
\begin{lstlisting}
struct RangeFenwick {
  int n;
  vector<int> b1, b2;
  RangeFenwick(int n) : n(n), b1(n + 2, 0), b2(n + 2, 0) {}
  void add_internal(vector<int>& b, int i, int v) {
    for (; i <= n; i += i & -i) b[i] += v;
  }
  void range_add(int l, int r, int v) {
    add_internal(b1, l, v); add_internal(b1, r + 1, -v);
    add_internal(b2, l, v * (l - 1)); add_internal(b2, r + 1, -v * r);
  }
  int prefix_sum(int i) {
    int s1 = 0, s2 = 0;
    for (int x = i; x > 0; x -= x & -x) s1 += b1[x], s2 += b2[x];
    return s1 * i - s2;
  }
  int range_sum(int l, int r) { return prefix_sum(r) - prefix_sum(l - 1); }
};
\end{lstlisting}

\algotitle{5.6 2D Fenwick Tree}
\algodesc{Point update and 2D prefix/submatrix sum in $O(\log N \log M)$.}
\begin{lstlisting}
struct Fenwick2D {
  int n, m;
  vector<vector<int>> tree;
  Fenwick2D(int n, int m) : n(n), m(m), tree(n + 1, vector<int>(m + 1, 0)) {}
  void add(int x, int y, int v) {
    for (int i = x; i <= n; i += i & -i)
      for (int j = y; j <= m; j += j & -j) tree[i][j] += v;
  }
  int query(int x, int y) {
    int sum = 0;
    for (int i = x; i > 0; i -= i & -i)
      for (int j = y; j > 0; j -= j & -j) sum += tree[i][j];
    return sum;
  }
  int query(int x1, int y1, int x2, int y2) {
    return query(x2, y2) - query(x1 - 1, y2) - query(x2, y1 - 1) + query(x1 - 1, y1 - 1);
  }
};
\end{lstlisting}

\algotitle{5.7 Segment Tree (User's Point Update, Range Sum)}
\algodesc{Zero-indexed array implementation matching user snippet.}
\begin{lstlisting}
class SGTree {
public:
  vector<int> seg;
  SGTree(int n) { seg.resize(4 * n + 1, 0); }
  void build(int idx, int l, int r, vector<int> &v) {
    if (l == r) { seg[idx] = v[l]; return; }
    int mid = (l + r) / 2;
    build(2 * idx + 1, l, mid, v);
    build(2 * idx + 2, mid + 1, r, v);
    seg[idx] = seg[2 * idx + 1] + seg[2 * idx + 2];
  }
  void update(int idx, int l, int r, int pos, int val) {
    if (l == r) { seg[idx] = val; return; }
    int mid = (l + r) / 2;
    if (pos <= mid) update(2 * idx + 1, l, mid, pos, val);
    else update(2 * idx + 2, mid + 1, r, pos, val);
    seg[idx] = seg[2 * idx + 1] + seg[2 * idx + 2];
  }
  int query(int idx, int l, int r, int ql, int qr) {
    if (qr < l || r < ql) return 0;
    if (ql <= l && r <= qr) return seg[idx];
    int mid = (l + r) / 2;
    return query(2 * idx + 1, l, mid, ql, qr) + query(2 * idx + 2, mid + 1, r, ql, qr);
  }
};
\end{lstlisting}

\algotitle{5.8 Lazy Segment Tree (Range Add \& Range Sum/Min/Max)}
\algodesc{General range update and range query in $O(\log N)$.}
\begin{lstlisting}
struct LazySeg {
  int n;
  vector<int> tree, lazy;
  LazySeg(int n) : n(n), tree(4 * n + 1, 0), lazy(4 * n + 1, 0) {}
  void push(int node, int l, int r) {
    if (!lazy[node]) return;
    tree[node] += (r - l + 1) * lazy[node];
    if (l != r) {
      lazy[node * 2] += lazy[node];
      lazy[node * 2 + 1] += lazy[node];
    }
    lazy[node] = 0;
  }
  void update(int node, int l, int r, int ql, int qr, int val) {
    push(node, l, r);
    if (qr < l || r < ql) return;
    if (ql <= l && r <= qr) {
      lazy[node] += val;
      push(node, l, r);
      return;
    }
    int mid = (l + r) / 2;
    update(node * 2, l, mid, ql, qr, val);
    update(node * 2 + 1, mid + 1, r, ql, qr, val);
    tree[node] = tree[node * 2] + tree[node * 2 + 1];
  }
  int query(int node, int l, int r, int ql, int qr) {
    push(node, l, r);
    if (qr < l || r < ql) return 0;
    if (ql <= l && r <= qr) return tree[node];
    int mid = (l + r) / 2;
    return query(node * 2, l, mid, ql, qr) + query(node * 2 + 1, mid + 1, r, ql, qr);
  }
};
\end{lstlisting}

\algotitle{5.9 Persistent Segment Tree (K-th Smallest in Range)}
\algodesc{Creates new nodes on update. Queries $k$-th element in range $[l, r]$ in $O(\log N)$.}
\begin{lstlisting}
struct PST {
  static const int MAXN = 200005, LG = 20;
  int lc[MAXN * LG], rc[MAXN * LG], sum[MAXN * LG], roots[MAXN], sz = 1;
  int update(int prev, int l, int r, int pos, int val) {
    int cur = sz++;
    lc[cur] = lc[prev]; rc[cur] = rc[prev]; sum[cur] = sum[prev] + val;
    if (l == r) return cur;
    int mid = (l + r) / 2;
    if (pos <= mid) lc[cur] = update(lc[prev], l, mid, pos, val);
    else rc[cur] = update(rc[prev], mid + 1, r, pos, val);
    return cur;
  }
  int query_kth(int u, int v, int l, int r, int k) {
    if (l == r) return l;
    int count = sum[lc[v]] - sum[lc[u]];
    int mid = (l + r) / 2;
    if (k <= count) return query_kth(lc[u], lc[v], l, mid, k);
    return query_kth(rc[u], rc[v], mid + 1, r, k - count);
  }
};
\end{lstlisting}

\algotitle{5.10 Sparse Table (1D RMQ in $O(1)$)}
\algodesc{Idempotent range minimum / maximum / gcd queries. Precomputation $O(N \log N)$.}
\begin{lstlisting}
struct SparseTable {
  int n, K;
  vector<vector<int>> st;
  SparseTable(const vector<int>& a) : n(a.size()), K(__lg(n) + 1) {
    st.assign(K, vector<int>(n));
    st[0] = a;
    for (int j = 1; j < K; j++)
      for (int i = 0; i + (1 << j) <= n; i++)
        st[j][i] = min(st[j - 1][i], st[j - 1][i + (1 << (j - 1))]);
  }
  int query(int l, int r) {
    int j = __lg(r - l + 1);
    return min(st[j][l], st[j][r - (1 << j) + 1]);
  }
};
\end{lstlisting}

\algotitle{5.11 Disjoint Sparse Table (Non-Idempotent $O(1)$)}
\algodesc{Answers arbitrary associative range queries (sum, XOR, matrix mul) in $O(1)$.}
\begin{lstlisting}
struct DisjointSparseTable {
  int n, levels;
  vector<vector<int>> table;
  DisjointSparseTable(vector<int> a) {
    n = a.size();
    int m = 1;
    while (m < n) m <<= 1;
    a.resize(m, 0);
    levels = __lg(m) + 1;
    table.assign(levels, vector<int>(m));
    table[0] = a;
    for (int h = 1; h < levels; h++) {
      int half = 1 << (h - 1);
      for (int i = 0; i < m; i += half * 2) {
        int mid = i + half;
        table[h][mid - 1] = a[mid - 1];
        for (int j = mid - 2; j >= i; j--) table[h][j] = table[h][j + 1] + a[j];
        table[h][mid] = a[mid];
        for (int j = mid + 1; j < i + half * 2; j++) table[h][j] = table[h][j - 1] + a[j];
      }
    }
  }
  int query(int l, int r) {
    if (l == r) return table[0][l];
    int h = __lg(l ^ r) + 1;
    return table[h][l] + table[h][r];
  }
};
\end{lstlisting}

\algotitle{5.12 Alphabet Trie \& Binary Trie (Max XOR)}
\algodesc{Prefix trees for strings and numbers. Max XOR pair query in $O(B)$ time.}
\begin{lstlisting}
// String Trie:
struct TrieNode {
  TrieNode* nxt[26] = {};
  bool is_end = false;
};
void insert_str(TrieNode* root, const string &s) {
  TrieNode* cur = root;
  for (char c : s) {
    int x = c - 'a';
    if (!cur->nxt[x]) cur->nxt[x] = new TrieNode();
    cur = cur->nxt[x];
  }
  cur->is_end = true;
}

// Binary Trie for Max XOR:
struct BitTrie {
  struct Node { int nxt[2] = {-1, -1}; int cnt = 0; };
  vector<Node> tree;
  BitTrie() { tree.push_back({}); }
  void insert(int x) {
    int cur = 0; tree[cur].cnt++;
    for (int i = 30; i >= 0; i--) {
      int b = (x >> i) & 1;
      if (tree[cur].nxt[b] == -1) {
        tree[cur].nxt[b] = tree.size();
        tree.push_back({});
      }
      cur = tree[cur].nxt[b];
      tree[cur].cnt++;
    }
  }
  int max_xor(int x) {
    int cur = 0, res = 0;
    for (int i = 30; i >= 0; i--) {
      int b = (x >> i) & 1;
      if (tree[cur].nxt[b ^ 1] != -1 && tree[tree[cur].nxt[b ^ 1]].cnt > 0) {
        res |= (1 << i);
        cur = tree[cur].nxt[b ^ 1];
      } else {
        cur = tree[cur].nxt[b];
      }
    }
    return res;
  }
};
\end{lstlisting}

\algotitle{5.13 Monotonic Stack \& Monotonic Queue}
\algodesc{Nearest Smaller to Left in $O(N)$; Sliding window maximum in $O(1)$ amortized.}
\begin{lstlisting}
// Nearest Smaller to Left (Values):
vector<int> nearest_smaller_left(const vector<int>& a) {
  int n = a.size();
  vector<int> ans(n, -1);
  stack<int> st;
  for (int i = 0; i < n; i++) {
    while (!st.empty() && st.top() >= a[i]) st.pop();
    if (!st.empty()) ans[i] = st.top();
    st.push(a[i]);
  }
  return ans;
}

// Sliding Window Maximum (Window size K):
vector<int> sliding_max(const vector<int>& a, int k) {
  deque<int> dq;
  vector<int> res;
  for (int i = 0; i < (int)a.size(); i++) {
    while (!dq.empty() && dq.front() <= i - k) dq.pop_front();
    while (!dq.empty() && a[dq.back()] <= a[i]) dq.pop_back();
    dq.push_back(i);
    if (i >= k - 1) res.push_back(a[dq.front()]);
  }
  return res;
}
\end{lstlisting}

\algotitle{5.14 Mo's Algorithm (Snake Order)}
\algodesc{Offline range queries in $O((N + Q)\sqrt{N})$. Alternating sorting cuts pointer moves.}
\begin{lstlisting}
struct Query {
  int l, r, id;
};
void solve_mo(int n, vector<Query>& queries) {
  int B = max(1LL, (int)sqrt(n));
  sort(queries.begin(), queries.end(), [&](const Query& a, const Query& b) {
    if (a.l / B != b.l / B) return a.l / B < b.l / B;
    return (a.l / B) & 1 ? a.r < b.r : a.r > b.r;
  });
  int curL = 0, curR = -1;
  auto add = [&](int idx) { /* include a[idx] */ };
  auto remove = [&](int idx) { /* remove a[idx] */ };
  for (auto &q : queries) {
    while (curL > q.l) add(--curL);
    while (curR < q.r) add(++curR);
    while (curL < q.l) remove(curL++);
    while (curR > q.r) remove(curR--);
    // ans[q.id] = current_answer;
  }
}
\end{lstlisting}

\algotitle{5.15 Small-to-Large Merging}
\algodesc{When merging data sets between components, always merge the smaller into the larger for $O(N \log N)$ total time.}
\begin{lstlisting}
void merge_sets(vector<int>& comp_a, vector<int>& comp_b) {
  if (comp_a.size() < comp_b.size()) swap(comp_a, comp_b);
  for (int val : comp_b) comp_a.push_back(val);
  comp_b.clear();
}
\end{lstlisting}

\algotitle{5.16 Coordinate Compression}
\algodesc{Maps large coordinate values to $[0, K-1]$ while preserving order.}
\begin{lstlisting}
vector<int> vals = a;
sort(vals.begin(), vals.end());
vals.erase(unique(vals.begin(), vals.end()), vals.end());
for (int &x : a) x = lower_bound(vals.begin(), vals.end(), x) - vals.begin();
\end{lstlisting}



\catsection{6. String Algorithms}

\algotitle{6.1 Double Polynomial Rolling Hash}
\algodesc{Static $O(1)$ substring hash queries with two moduli to eliminate collision risks.}
\begin{lstlisting}
struct StringHash {
  int n;
  const int mod1 = 1000000007, mod2 = 1000000009;
  const int base1 = 313, base2 = 317;
  vector<int> h1, h2, p1, p2;
  StringHash(const string& s) : n(s.size()), h1(n + 1, 0), h2(n + 1, 0), p1(n + 1, 1), p2(n + 1, 1) {
    for (int i = 0; i < n; i++) {
      p1[i + 1] = (__int128)p1[i] * base1 % mod1;
      p2[i + 1] = (__int128)p2[i] * base2 % mod2;
      h1[i + 1] = ((__int128)h1[i] * base1 + s[i]) % mod1;
      h2[i + 1] = ((__int128)h2[i] * base2 + s[i]) % mod2;
    }
  }
  pair<int, int> get_hash(int l, int r) { // 0-indexed [l, r]
    int x1 = (h1[r + 1] - (__int128)h1[l] * p1[r - l + 1]) % mod1;
    if (x1 < 0) x1 += mod1;
    int x2 = (h2[r + 1] - (__int128)h2[l] * p2[r - l + 1]) % mod2;
    if (x2 < 0) x2 += mod2;
    return {x1, x2};
  }
};
\end{lstlisting}

\algotitle{6.2 Dynamic String Hash (Point Update Segment Tree)}
\algodesc{Supports single-character replacement and $O(\log N)$ substring hash queries.}
\begin{lstlisting}
struct DynamicHash {
  int n; string s;
  vector<int> h, p;
  const int base = 313, mod = 1000000007;
  DynamicHash(string str) : n(str.size()), s(str), h(4 * n + 1), p(n + 1, 1) {
    for (int i = 1; i <= n; i++) p[i] = (int)((__int128)p[i - 1] * base % mod);
    build(1, 0, n - 1);
  }
  void build(int node, int l, int r) {
    if (l == r) { h[node] = s[l]; return; }
    int mid = (l + r) / 2;
    build(node * 2, l, mid);
    build(node * 2 + 1, mid + 1, r);
    h[node] = ((__int128)h[node * 2] * p[r - mid] + h[node * 2 + 1]) % mod;
  }
  void update(int node, int l, int r, int pos, char c) {
    if (l == r) { h[node] = s[pos] = c; return; }
    int mid = (l + r) / 2;
    if (pos <= mid) update(node * 2, l, mid, pos, c);
    else update(node * 2 + 1, mid + 1, r, pos, c);
    h[node] = ((__int128)h[node * 2] * p[r - mid] + h[node * 2 + 1]) % mod;
  }
  int query(int node, int l, int r, int ql, int qr) {
    if (ql <= l && r <= qr) return h[node];
    int mid = (l + r) / 2;
    if (qr <= mid) return query(node * 2, l, mid, ql, qr);
    if (ql > mid) return query(node * 2 + 1, mid + 1, r, ql, qr);
    int left_h = query(node * 2, l, mid, ql, qr);
    int right_h = query(node * 2 + 1, mid + 1, r, ql, qr);
    return ((__int128)left_h * p[min(r, qr) - mid] + right_h) % mod;
  }
};
\end{lstlisting}

\algotitle{6.3 Prefix Function (KMP) \& Pattern Matching}
\algodesc{$\pi[i]$ is the length of the longest proper prefix of $s[0..i]$ that is also a suffix of $s[0..i]$.}
\begin{lstlisting}
vector<int> prefix_function(const string& s) {
  int n = s.size();
  vector<int> pi(n, 0);
  for (int i = 1; i < n; i++) {
    int j = pi[i - 1];
    while (j > 0 && s[i] != s[j]) j = pi[j - 1];
    if (s[i] == s[j]) j++;
    pi[i] = j;
  }
  return pi;
}
vector<int> kmp_search(const string& text, const string& pat) {
  string cur = pat + "#" + text;
  vector<int> pi = prefix_function(cur), matches;
  for (int i = pat.size() + 1; i < (int)cur.size(); i++) {
    if (pi[i] == (int)pat.size()) matches.push_back(i - 2 * pat.size());
  }
  return matches;
}
\end{lstlisting}

\algotitle{6.4 Z-Algorithm}
\algodesc{$Z[i]$ is the length of the longest common prefix of $s$ and $s[i..n-1]$ in $O(N)$ time.}
\begin{lstlisting}
vector<int> z_function(const string& s) {
  int n = s.size();
  vector<int> z(n, 0);
  for (int i = 1, l = 0, r = 0; i < n; i++) {
    if (i <= r) z[i] = min(r - i + 1, z[i - l]);
    while (i + z[i] < n && s[z[i]] == s[i + z[i]]) z[i]++;
    if (i + z[i] - 1 > r) { l = i; r = i + z[i] - 1; }
  }
  return z;
}
\end{lstlisting}

\algotitle{6.5 Manacher's Algorithm (Palindromes in $O(N)$)}
\algodesc{Computes all palindrome radii in $O(N)$. Interleaving with `\#' handles odd and even palindromes uniformly.}
\begin{itemize}[leftmargin=*,noitemsep,topsep=1pt]
  \item Simple Explanation: Transform string ``aba'' to ``@\#a\#b\#a\#\$''. Then every palindrome (odd or even) has an odd length with a single center.
  \item Radius Reuse: When expanding around center $i$, if $i < r$, we copy the previously computed radius from its symmetric mirror $i' = l + r - i$: $p[i] = \min(r - i, p[i'])$. We only expand character-by-character if the palindrome extends beyond the current known rightmost boundary $r$. Since $r$ strictly increases, the overall time is $O(N)$.
\end{itemize}
\begin{lstlisting}
// Returns array p where p[i] is radius in transformed string
vector<int> manacher(const string& s) {
  string t = "@";
  for (char c : s) { t += "#"; t += c; }
  t += "#$";
  int n = t.size(), l = 1, r = 1;
  vector<int> p(n, 0);
  for (int i = 1; i < n - 1; i++) {
    p[i] = max(0LL, min(r - i, p[l + r - i]));
    while (t[i - p[i]] == t[i + p[i]]) p[i]++;
    if (i + p[i] > r) { l = i - p[i]; r = i + p[i]; }
  }
  return p; // Longest palindrome substring length is max_element(p) - 1
}
bool is_palindrome(const vector<int>& p, int orig_l, int orig_r) {
  int center = orig_l + orig_r + 2; // In 0-indexed original string
  int len = orig_r - orig_l + 1;
  return p[center] - 1 >= len;
}
\end{lstlisting}

\algotitle{6.6 Suffix Array \& LCP ($O(N \log N)$)}
\algodesc{Constructs suffix array using prefix doubling. Kasai's algorithm finds LCP array in $O(N)$.}
\begin{lstlisting}
struct SuffixArray {
  int n;
  vector<int> sa, lcp, rank;
  SuffixArray(string s) : n(s.size() + 1), sa(n), lcp(n), rank(n) {
    s += "$";
    vector<int> k(n);
    for (int i = 0; i < n; i++) { sa[i] = i; rank[i] = s[i]; }
    for (int len = 1; len < n; len <<= 1) {
      auto cmp = [&](int i, int j) {
        if (rank[i] != rank[j]) return rank[i] < rank[j];
        int ri = (i + len < n ? rank[i + len] : -1);
        int rj = (j + len < n ? rank[j + len] : -1);
        return ri < rj;
      };
      sort(sa.begin(), sa.end(), cmp);
      k[sa[0]] = 0;
      for (int i = 1; i < n; i++) k[sa[i]] = k[sa[i - 1]] + cmp(sa[i - 1], sa[i]);
      rank = k;
      if (rank[sa[n - 1]] == n - 1) break;
    }
    // Kasai's algorithm for LCP:
    for (int i = 0, h = 0; i < n - 1; i++) {
      int pi = rank[i], j = sa[pi - 1];
      while (s[i + h] == s[j + h]) h++;
      lcp[pi] = h;
      if (h > 0) h--;
    }
  }
  // Total distinct substrings: n*(n-1)/2 - sum(lcp[i])
};
\end{lstlisting}

\algotitle{6.7 Aho-Corasick Automaton}
\algodesc{Multi-pattern string matching in $O(\text{Text} + \sum |\text{Patterns}|)$. Builds BFS failure links.}
\begin{lstlisting}
struct AhoCorasick {
  struct Node { int nxt[26] = {}, fail = 0, match_cnt = 0; };
  vector<Node> t;
  AhoCorasick() { t.push_back({}); }
  void insert(const string& s) {
    int cur = 0;
    for (char c : s) {
      int x = c - 'a';
      if (!t[cur].nxt[x]) { t[cur].nxt[x] = t.size(); t.push_back({}); }
      cur = t[cur].nxt[x];
    }
    t[cur].match_cnt++;
  }
  void build() {
    queue<int> q;
    for (int c = 0; c < 26; c++) if (t[0].nxt[c]) q.push(t[0].nxt[c]);
    while (!q.empty()) {
      int u = q.front(); q.pop();
      for (int c = 0; c < 26; c++) {
        if (t[u].nxt[c]) {
          t[t[u].nxt[c]].fail = t[t[u].fail].nxt[c];
          t[t[u].nxt[c]].match_cnt += t[t[t[u].nxt[c]].fail].match_cnt;
          q.push(t[u].nxt[c]);
        } else {
          t[u].nxt[c] = t[t[u].fail].nxt[c];
        }
      }
    }
  }
};
\end{lstlisting}

\algotitle{6.8 Booth's Algorithm (Minimum Lexicographical Rotation)}
\algodesc{Finds starting index of lexicographically smallest cyclic shift of $s$ in $O(N)$ time.}
\begin{lstlisting}
int min_cyclic_string(string s) {
  s += s;
  int n = s.size(), i = 0, ans = 0;
  while (i < n / 2) {
    ans = i;
    int j = i + 1, k = i;
    while (j < n && s[k] <= s[j]) {
      if (s[k] < s[j]) k = i;
      else k++;
      j++;
    }
    while (i <= k) i += j - k;
  }
  return ans;
}
\end{lstlisting}

\algotitle{6.9 String Utility Helpers}
\algodesc{Fast conversions, tokenization, and character properties.}
\begin{lstlisting}
// Tokenize by whitespace:
vector<string> split(const string& s) {
  stringstream ss(s); string w;
  vector<string> res;
  while (ss >> w) res.push_back(w);
  return res;
}
// String <-> Number: stoi(s), stoll(s), stod(s), to_string(x)
// Char tests: islower(c), isupper(c), isdigit(c), isalpha(c), tolower(c), toupper(c)
\end{lstlisting}



\catsection{7. Graph Theory}

\algotitle{7.1 Graph Traversals (DFS, BFS, 0-1 BFS, Multi-Source)}
\algodesc{Core search algorithms for unweighted and 0/1 weighted graphs.}
\begin{lstlisting}
// 0-1 BFS in O(V + E) using std::deque:
vector<int> bfs01(int src, int n, const vector<vector<pair<int, int>>>& g) {
  vector<int> dist(n + 1, INF);
  deque<int> dq;
  dist[src] = 0; dq.push_front(src);
  while (!dq.empty()) {
    int u = dq.front(); dq.pop_front();
    for (auto [v, w] : g[u]) {
      if (dist[v] > dist[u] + w) {
        dist[v] = dist[u] + w;
        if (w == 0) dq.push_front(v);
        else dq.push_back(v);
      }
    }
  }
  return dist;
}
\end{lstlisting}

\algotitle{7.2 Dijkstra's Algorithm (Single Source Shortest Path)}
\algodesc{Finds shortest paths with non-negative edge weights in $O((V + E)\log V)$.}
\begin{lstlisting}
vector<int> dijkstra(int s, int n, const vector<vector<pair<int, int>>>& g) {
  vector<int> dist(n + 1, INF);
  priority_queue<pair<int, int>, vector<pair<int, int>>, greater<pair<int, int>>> pq;
  dist[s] = 0; pq.push({0, s});
  while (!pq.empty()) {
    auto [d, u] = pq.top(); pq.pop();
    if (d > dist[u]) continue;
    for (auto [v, w] : g[u]) {
      if (dist[v] > dist[u] + w) {
        dist[v] = dist[u] + w;
        pq.push({dist[v], v});
      }
    }
  }
  return dist;
}
\end{lstlisting}

\algotitle{7.3 Bellman-Ford \& Negative Cycle Detection}
\algodesc{Finds shortest paths with negative edges in $O(V \cdot E)$ and identifies negative cycles.}
\begin{lstlisting}
struct Edge { int u, v, w; };
vector<int> bellman_ford(int s, int n, const vector<Edge>& edges, bool &has_neg_cycle) {
  vector<int> dist(n + 1, INF);
  dist[s] = 0;
  for (int i = 1; i <= n - 1; i++) {
    for (const auto& e : edges) {
      if (dist[e.u] < INF && dist[e.v] > dist[e.u] + e.w) dist[e.v] = dist[e.u] + e.w;
    }
  }
  has_neg_cycle = false;
  for (const auto& e : edges) {
    if (dist[e.u] < INF && dist[e.v] > dist[e.u] + e.w) {
      has_neg_cycle = true;
      dist[e.v] = -INF; // Reachable from negative cycle
    }
  }
  return dist;
}
\end{lstlisting}

\algotitle{7.4 Floyd-Warshall (All-Pairs Shortest Path)}
\algodesc{Computes all-pairs shortest paths in $O(V^3)$. Propagates negative cycles.}
\begin{lstlisting}
void floyd_warshall(int n, vector<vector<int>>& d) {
  for (int i = 1; i <= n; i++) d[i][i] = min(d[i][i], 0LL);
  for (int k = 1; k <= n; k++)
    for (int i = 1; i <= n; i++)
      for (int j = 1; j <= n; j++)
        if (d[i][k] < INF && d[k][j] < INF)
          d[i][j] = min(d[i][j], d[i][k] + d[k][j]);
  // Propagate -INF for negative cycles:
  for (int k = 1; k <= n; k++)
    if (d[k][k] < 0)
      for (int i = 1; i <= n; i++)
        for (int j = 1; j <= n; j++)
          if (d[i][k] < INF && d[k][j] < INF) d[i][j] = -INF;
}
\end{lstlisting}

\algotitle{7.5 Strongly Connected Components (Kosaraju's)}
\algodesc{Two DFS passes: 1st on $G$ for finish times; 2nd on $G^T$ to extract SCCs. Builds condensation DAG.}
\begin{lstlisting}
struct KosarajuSCC {
  int n, comps = 0;
  vector<vector<int>> g, rg, dag;
  vector<int> order, comp;
  vector<bool> vis;
  KosarajuSCC(int n) : n(n), g(n + 1), rg(n + 1), comp(n + 1, -1), vis(n + 1, false) {}
  void add_edge(int u, int v) { g[u].push_back(v); rg[v].push_back(u); }
  void dfs1(int u) {
    vis[u] = true;
    for (int v : g[u]) if (!vis[v]) dfs1(v);
    order.push_back(u);
  }
  void dfs2(int u, int c) {
    comp[u] = c;
    for (int v : rg[u]) if (comp[v] == -1) dfs2(v, c);
  }
  void build() {
    for (int i = 1; i <= n; i++) if (!vis[i]) dfs1(i);
    reverse(order.begin(), order.end());
    for (int u : order) if (comp[u] == -1) dfs2(u, comps++);
    dag.resize(comps);
    for (int u = 1; u <= n; u++) {
      for (int v : g[u]) if (comp[u] != comp[v]) dag[comp[u]].push_back(comp[v]);
    }
  }
};
\end{lstlisting}

\algotitle{7.6 Strongly Connected Components (Tarjan's)}
\algodesc{Single-pass DFS with explicit stack and entry/low values to find SCCs in $O(V + E)$.}
\begin{lstlisting}
struct TarjanSCC {
  int n, timer = 0, comps = 0;
  vector<vector<int>> g;
  vector<int> tin, low, comp;
  vector<bool> in_st;
  stack<int> st;
  TarjanSCC(int n) : n(n), g(n + 1), tin(n + 1, -1), low(n + 1, -1), comp(n + 1, -1), in_st(n + 1, false) {}
  void add_edge(int u, int v) { g[u].push_back(v); }
  void dfs(int u) {
    tin[u] = low[u] = ++timer;
    st.push(u); in_st[u] = true;
    for (int v : g[u]) {
      if (tin[v] == -1) {
        dfs(v);
        low[u] = min(low[u], low[v]);
      } else if (in_st[v]) {
        low[u] = min(low[u], tin[v]);
      }
    }
    if (low[u] == tin[u]) {
      while (true) {
        int v = st.top(); st.pop();
        in_st[v] = false;
        comp[v] = comps;
        if (v == u) break;
      }
      comps++;
    }
  }
  void build() { for (int i = 1; i <= n; i++) if (tin[i] == -1) dfs(i); }
};
\end{lstlisting}

\algotitle{7.7 2-SAT Solver (via SCC)}
\algodesc{Solves 2-SAT clause constraints $(u \lor v)$. Assigns true/false values in $O(V + E)$.}
\begin{lstlisting}
struct TwoSAT {
  int n;
  KosarajuSCC scc;
  TwoSAT(int n) : n(n), scc(2 * n) {}
  // Node x: 2*x (false), 2*x + 1 (true)
  void add_clause(int u, bool val_u, int v, bool val_v) {
    int pu = 2 * u + (val_u ? 1 : 0);
    int pv = 2 * v + (val_v ? 1 : 0);
    scc.add_edge(pu ^ 1, pv); // ~u => v
    scc.add_edge(pv ^ 1, pu); // ~v => u
  }
  bool solve(vector<bool>& assignment) {
    scc.build();
    assignment.resize(n + 1);
    for (int i = 1; i <= n; i++) {
      if (scc.comp[2 * i] == scc.comp[2 * i + 1]) return false; // Unsatisfiable
      assignment[i] = scc.comp[2 * i + 1] < scc.comp[2 * i];
    }
    return true;
  }
};
\end{lstlisting}

\algotitle{7.8 Articulation Points \& Bridges}
\algodesc{Finds all cut vertices and bridge edges in an undirected graph in $O(V + E)$.}
\begin{lstlisting}
struct BridgesAndArt {
  int n, timer = 0;
  vector<vector<int>> g;
  vector<int> tin, low;
  vector<bool> is_art;
  vector<pair<int, int>> bridges;
  BridgesAndArt(int n) : n(n), g(n + 1), tin(n + 1, -1), low(n + 1, -1), is_art(n + 1, false) {}
  void add_edge(int u, int v) { g[u].push_back(v); g[v].push_back(u); }
  void dfs(int u, int p = -1) {
    tin[u] = low[u] = ++timer;
    int children = 0;
    for (int v : g[u]) {
      if (v == p) continue;
      if (tin[v] != -1) {
        low[u] = min(low[u], tin[v]);
      } else {
        dfs(v, u);
        low[u] = min(low[u], low[v]);
        if (low[v] > tin[u]) bridges.push_back({u, v});
        if (low[v] >= tin[u] && p != -1) is_art[u] = true;
        children++;
      }
    }
    if (p == -1 && children > 1) is_art[u] = true;
  }
  void build() { for (int i = 1; i <= n; i++) if (tin[i] == -1) dfs(i); }
};
\end{lstlisting}

\algotitle{7.9 Tree Flattening (Euler Tour)}
\algodesc{Maps tree subtrees into contiguous 1D array segments $[tin[u], tout[u]]$ for range data structures.}
\begin{lstlisting}
vector<int> tin, tout;
int timer = 0;
void euler_tour(int u, int p, const vector<vector<int>>& g) {
  tin[u] = ++timer;
  for (int v : g[u]) if (v != p) euler_tour(v, u, g);
  tout[u] = timer;
}
\end{lstlisting}

\algotitle{7.10 Binary Lifting \& LCA}
\algodesc{Precomputation $O(N \log N)$; LCA and $k$-th ancestor queries in $O(\log N)$.}
\begin{lstlisting}
struct LCA {
  int n, LG;
  vector<vector<int>> up;
  vector<int> depth;
  LCA(int n, const vector<vector<int>>& g, int root = 1) : n(n), LG(__lg(n) + 1), up(LG, vector<int>(n + 1, 0)), depth(n + 1, 0) {
    auto dfs = [&](auto self, int u, int p, int d) -> void {
      up[0][u] = p; depth[u] = d;
      for (int i = 1; i < LG; i++) up[i][u] = up[i - 1][up[i - 1][u]];
      for (int v : g[u]) if (v != p) self(self, v, u, d + 1);
    };
    dfs(dfs, root, 0, 0);
  }
  int get_lca(int u, int v) {
    if (depth[u] < depth[v]) swap(u, v);
    for (int i = LG - 1; i >= 0; i--)
      if (depth[u] - (1 << i) >= depth[v]) u = up[i][u];
    if (u == v) return u;
    for (int i = LG - 1; i >= 0; i--) {
      if (up[i][u] != up[i][v]) { u = up[i][u]; v = up[i][v]; }
    }
    return up[0][u];
  }
};
\end{lstlisting}

\algotitle{7.11 Heavy-Light Decomposition (HLD)}
\algodesc{Decomposes tree into $O(\log N)$ heavy paths. Supports path queries and updates via segment tree.}
\begin{lstlisting}
struct HLD {
  int n, cur_pos = 1;
  vector<int> parent, depth, heavy, head, pos, sz;
  vector<vector<int>> g;
  HLD(int n) : n(n), parent(n + 1), depth(n + 1), heavy(n + 1, -1), head(n + 1), pos(n + 1), sz(n + 1), g(n + 1) {}
  void add_edge(int u, int v) { g[u].push_back(v); g[v].push_back(u); }
  int dfs(int u, int p, int d) {
    parent[u] = p; depth[u] = d; sz[u] = 1;
    int max_c = 0;
    for (int v : g[u]) {
      if (v == p) continue;
      int c_sz = dfs(v, u, d + 1);
      sz[u] += c_sz;
      if (c_sz > max_c) { max_c = c_sz; heavy[u] = v; }
    }
    return sz[u];
  }
  void decompose(int u, int h) {
    head[u] = h; pos[u] = cur_pos++;
    if (heavy[u] != -1) decompose(heavy[u], h);
    for (int v : g[u]) {
      if (v != parent[u] && v != heavy[u]) decompose(v, v);
    }
  }
  void init(int root = 1) { dfs(root, 0, 0); decompose(root, root); }
};
\end{lstlisting}

\algotitle{7.12 DSU on Tree (Sack / Small-to-Large on Trees)}
\algodesc{Solves offline subtree queries in $O(N \log N)$ by keeping the heavy child's contribution.}
\begin{lstlisting}
int sz_tree[N], heavy_child[N];
void dfs_sz(int u, int p, const vector<vector<int>>& g) {
  sz_tree[u] = 1; heavy_child[u] = -1; int max_c = 0;
  for (int v : g[u]) if (v != p) {
    dfs_sz(v, u, g);
    sz_tree[u] += sz_tree[v];
    if (sz_tree[v] > max_c) { max_c = sz_tree[v]; heavy_child[u] = v; }
  }
}
void dfs_sack(int u, int p, bool keep, const vector<vector<int>>& g) {
  for (int v : g[u]) if (v != p && v != heavy_child[u]) dfs_sack(v, u, false, g);
  if (heavy_child[u] != -1) dfs_sack(heavy_child[u], u, true, g);
  // Add light children subtrees and node u to DS
  // Answer queries for u
  if (!keep) { /* Remove all elements in subtree u from DS */ }
}
\end{lstlisting}

\algotitle{7.13 Dinic's Algorithm (Maximum Flow)}
\algodesc{Finds maximum flow in $O(V^2 E)$ (or $O(E\sqrt{V})$ for unit networks) using BFS levels and DFS blocking flow.}
\begin{lstlisting}
struct Dinic {
  struct Edge { int to, rev; int cap, flow; };
  int n, s, t;
  vector<vector<Edge>> adj;
  vector<int> level, ptr;
  Dinic(int n, int s, int t) : n(n), s(s), t(t), adj(n + 1), level(n + 1), ptr(n + 1) {}
  void add_edge(int u, int v, int cap) {
    adj[u].push_back({v, (int)adj[v].size(), cap, 0});
    adj[v].push_back({u, (int)adj[u].size() - 1, 0, 0});
  }
  bool bfs() {
    fill(level.begin(), level.end(), -1);
    level[s] = 0; queue<int> q; q.push(s);
    while (!q.empty()) {
      int u = q.front(); q.pop();
      for (const auto& e : adj[u]) {
        if (e.cap - e.flow > 0 && level[e.to] == -1) {
          level[e.to] = level[u] + 1; q.push(e.to);
        }
      }
    }
    return level[t] != -1;
  }
  int dfs(int u, int pushed) {
    if (pushed == 0 || u == t) return pushed;
    for (int& cid = ptr[u]; cid < (int)adj[u].size(); cid++) {
      auto& e = adj[u][cid];
      if (level[u] + 1 != level[e.to] || e.cap - e.flow == 0) continue;
      int tr = dfs(e.to, min(pushed, e.cap - e.flow));
      if (tr == 0) continue;
      e.flow += tr;
      adj[e.to][e.rev].flow -= tr;
      return tr;
    }
    return 0;
  }
  int max_flow() {
    int flow = 0;
    while (bfs()) {
      fill(ptr.begin(), ptr.end(), 0);
      while (int pushed = dfs(s, INF)) flow += pushed;
    }
    return flow;
  }
};
\end{lstlisting}

\algotitle{7.14 Maximum Bipartite Matching (Kuhn's)}
\algodesc{Simple $O(V \cdot E)$ augmenting path algorithm for bipartite matching.}
\begin{lstlisting}
bool dfs_kuhn(int u, const vector<vector<int>>& g, vector<int>& match, vector<bool>& vis) {
  if (vis[u]) return false;
  vis[u] = true;
  for (int v : g[u]) {
    if (match[v] == -1 || dfs_kuhn(match[v], g, match, vis)) {
      match[v] = u; return true;
    }
  }
  return false;
}
int max_bipartite_matching(int n_left, int n_right, const vector<vector<int>>& g) {
  vector<int> match(n_right + 1, -1);
  int res = 0;
  for (int u = 1; u <= n_left; u++) {
    vector<bool> vis(n_left + 1, false);
    if (dfs_kuhn(u, g, match, vis)) res++;
  }
  return res;
}
\end{lstlisting}

\algotitle{7.15 Tree Diameter (2-BFS)}
\algodesc{BFS from node 1 gives farthest node $u$; BFS from $u$ gives farthest node $v$ and diameter.}
\begin{lstlisting}
pair<int, int> bfs_farthest(int start, int n, const vector<vector<int>>& g) {
  vector<int> dist(n + 1, -1);
  queue<int> q; q.push(start); dist[start] = 0;
  int far_node = start;
  while (!q.empty()) {
    int u = q.front(); q.pop();
    if (dist[u] > dist[far_node]) far_node = u;
    for (int v : g[u]) if (dist[v] == -1) { dist[v] = dist[u] + 1; q.push(v); }
  }
  return {far_node, dist[far_node]};
}
int tree_diameter(int n, const vector<vector<int>>& g) {
  int u = bfs_farthest(1, n, g).first;
  return bfs_farthest(u, n, g).second;
}
\end{lstlisting}



\catsection{8. Dynamic Programming}

\algotitle{8.1 Classical DP (Knapsack, Coins, LIS, LCS, Edit Dist)}
\algodesc{Essential competitive programming dynamic programming templates.}
\begin{lstlisting}
// 0/1 Knapsack (Reverse iteration):
for (int i = 0; i < n; i++)
  for (int j = W; j >= wt[i]; j--) dp[j] = max(dp[j], dp[j - wt[i]] + val[i]);

// Unbounded Knapsack (Forward iteration):
for (int i = 0; i < n; i++)
  for (int j = wt[i]; j <= W; j++) dp[j] = max(dp[j], dp[j - wt[i]] + val[i]);

// LIS in O(N log N):
vector<int> tails;
for (int x : a) {
  auto it = lower_bound(tails.begin(), tails.end(), x);
  if (it == tails.end()) tails.push_back(x);
  else *it = x;
}
int lis_len = tails.size();

// LCS in O(N*M):
for (int i = 1; i <= n; i++)
  for (int j = 1; j <= m; j++)
    dp[i][j] = (s[i-1] == t[j-1] ? dp[i-1][j-1] + 1 : max(dp[i-1][j], dp[i][j-1]));

// Edit Distance:
for (int i = 1; i <= n; i++)
  for (int j = 1; j <= m; j++) {
    if (s[i-1] == t[j-1]) dp[i][j] = dp[i-1][j-1];
    else dp[i][j] = 1 + min({dp[i-1][j], dp[i][j-1], dp[i-1][j-1]});
  }
\end{lstlisting}

\algotitle{8.2 Kadane's Algorithm \& Max Submatrix Sum}
\algodesc{Max subarray sum in $O(N)$; 2D submatrix max sum in $O(N^2 M)$ by collapsing columns.}
\begin{lstlisting}
int kadane(const vector<int>& a) {
  int cur = 0, ans = -INF;
  for (int x : a) { cur = max(x, cur + x); ans = max(ans, cur); }
  return ans;
}
int max_submatrix_sum(int n, int m, const vector<vector<int>>& mat) {
  int best = -INF;
  for (int top = 0; top < n; top++) {
    vector<int> col_sum(m, 0);
    for (int bot = top; bot < n; bot++) {
      for (int j = 0; j < m; j++) col_sum[j] += mat[bot][j];
      best = max(best, kadane(col_sum));
    }
  }
  return best;
}
\end{lstlisting}

\algotitle{8.3 Subset Sum via Bitset (Binary Splitting)}
\algodesc{Determines reachable sums in $O(\frac{N\sqrt{N}}{64})$ by grouping identical weights into powers of 2.}
\begin{lstlisting}
bitset<1000005> reachable_subset_sums(const map<int, int>& item_counts) {
  bitset<1000005> bs;
  bs[0] = 1;
  for (auto [val, cnt] : item_counts) {
    for (int k = 1; k <= cnt; k <<= 1) {
      bs |= (bs << (val * k));
      cnt -= k;
    }
    if (cnt > 0) bs |= (bs << (val * cnt));
  }
  return bs;
}
\end{lstlisting}

\algotitle{8.4 Interval DP (Matrix Chain Multiplication)}
\algodesc{DP over range lengths: $dp[l][r] = \min_{l \le k < r} (dp[l][k] + dp[k+1][r] + \text{cost})$.}
\begin{lstlisting}
for (int len = 2; len <= n; len++) {
  for (int l = 1; l + len - 1 <= n; l++) {
    int r = l + len - 1;
    dp[l][r] = INF;
    for (int k = l; k < r; k++) {
      dp[l][r] = min(dp[l][r], dp[l][k] + dp[k + 1][r] + cost(l, k, r));
    }
  }
}
\end{lstlisting}

\algotitle{8.5 Digit DP Template}
\algodesc{Counts numbers in $[0, N]$ satisfying digit properties with memoization.}
\begin{lstlisting}
int dp[20][2][2][100]; // pos, tight, leading_zeros, state
string num_str;
int solve_digit(int pos, bool tight, bool lead, int state) {
  if (pos == (int)num_str.size()) return 1; // Valid number
  if (dp[pos][tight][lead][state] != -1) return dp[pos][tight][lead][state];
  int limit = tight ? (num_str[pos] - '0') : 9;
  int ans = 0;
  for (int dig = 0; dig <= limit; dig++) {
    bool next_tight = tight && (dig == limit);
    bool next_lead = lead && (dig == 0);
    int next_state = (next_lead ? 0 : state + dig);
    ans += solve_digit(pos + 1, next_tight, next_lead, next_state);
  }
  return dp[pos][tight][lead][state] = ans;
}
int count_up_to(int n) {
  num_str = to_string(n);
  memset(dp, -1, sizeof(dp));
  return solve_digit(0, true, true, 0);
}
\end{lstlisting}

\algotitle{8.6 Bitmask DP (TSP / Shortest Hamiltonian Cycle)}
\algodesc{Finds minimum weight Hamiltonian cycle visiting every vertex once in $O(2^N \cdot N^2)$.}
\begin{lstlisting}
int tsp(int n, const vector<vector<int>>& dist) {
  vector<vector<int>> dp(1 << n, vector<int>(n, INF));
  dp[1][0] = 0; // Start at vertex 0
  for (int mask = 1; mask < (1 << n); mask++) {
    for (int u = 0; u < n; u++) {
      if (dp[mask][u] == INF) continue;
      for (int v = 0; v < n; v++) {
        if (!(mask & (1 << v))) {
          int nxt = mask | (1 << v);
          dp[nxt][v] = min(dp[nxt][v], dp[mask][u] + dist[u][v]);
        }
      }
    }
  }
  int ans = INF;
  for (int u = 1; u < n; u++) ans = min(ans, dp[(1 << n) - 1][u] + dist[u][0]);
  return ans;
}
\end{lstlisting}

\algotitle{8.7 Tree DP (Maximum Weight Independent Set)}
\algodesc{Chooses a subset of nodes with no two adjacent so total weight is maximized.}
\begin{lstlisting}
int dp[N][2]; // dp[u][0]: u not picked; dp[u][1]: u picked
void dfs_tree_dp(int u, int p, const vector<vector<int>>& g, const vector<int>& weight) {
  dp[u][0] = 0;
  dp[u][1] = weight[u];
  for (int v : g[u]) {
    if (v == p) continue;
    dfs_tree_dp(v, u, g, weight);
    dp[u][0] += max(dp[v][0], dp[v][1]);
    dp[u][1] += dp[v][0];
  }
}
\end{lstlisting}

\algotitle{8.8 SOS DP (Sum Over Subsets / Fast Zeta Transform)}
\algodesc{Computes $F(mask) = \sum_{sub \subseteq mask} A(sub)$ in $O(B \cdot 2^B)$ time.}
\begin{lstlisting}
void sos_dp(vector<int>& f, int bits) {
  for (int b = 0; b < bits; b++) {
    for (int mask = 0; mask < (1 << bits); mask++) {
      if (mask & (1 << b)) f[mask] += f[mask ^ (1 << b)];
    }
  }
}
// Inverse Zeta Transform (Mobius transform on bitmasks):
void inverse_sos_dp(vector<int>& f, int bits) {
  for (int b = 0; b < bits; b++) {
    for (int mask = 0; mask < (1 << bits); mask++) {
      if (mask & (1 << b)) f[mask] -= f[mask ^ (1 << b)];
    }
  }
}
\end{lstlisting}

\algotitle{8.9 Knuth's DP Optimization}
\algodesc{Reduces $O(N^3)$ interval DP to $O(N^2)$ when optimal split point satisfies $opt[i][j-1] \le opt[i][j] \le opt[i+1][j]$.}
\begin{lstlisting}
void knuth_optimization(int n) {
  for (int i = 1; i <= n; i++) { opt[i][i] = i; dp[i][i] = 0; }
  for (int len = 2; len <= n; len++) {
    for (int i = 1; i + len - 1 <= n; i++) {
      int j = i + len - 1;
      dp[i][j] = INF;
      for (int k = opt[i][j - 1]; k <= min(j - 1, opt[i + 1][j]); k++) {
        int cost = dp[i][k] + dp[k + 1][j] + C(i, j);
        if (cost < dp[i][j]) { dp[i][j] = cost; opt[i][j] = k; }
      }
    }
  }
}
\end{lstlisting}

\algotitle{8.10 Li Chao Tree (Convex Hull Trick for Lines)}
\algodesc{Maintains set of lines $y = mx + b$; answers minimum/maximum $y$ at query $x$ in $O(\log C)$.}
\begin{lstlisting}
struct Line { int m, b; int eval(int x) { return m * x + b; } };
struct LiChaoTree {
  struct Node { Line line; int lc = -1, rc = -1; };
  vector<Node> tree;
  LiChaoTree() { tree.push_back({{0, INF}, -1, -1}); }
  void add_line(Line nw, int node, int l, int r) {
    int mid = (l + r) / 2;
    bool left = nw.eval(l) < tree[node].line.eval(l);
    bool m = nw.eval(mid) < tree[node].line.eval(mid);
    if (m) swap(tree[node].line, nw);
    if (l == r) return;
    if (left != m) {
      if (tree[node].lc == -1) { tree[node].lc = tree.size(); tree.push_back({{0, INF}, -1, -1}); }
      add_line(nw, tree[node].lc, l, mid);
    } else {
      if (tree[node].rc == -1) { tree[node].rc = tree.size(); tree.push_back({{0, INF}, -1, -1}); }
      add_line(nw, tree[node].rc, mid + 1, r);
    }
  }
  int query(int x, int node, int l, int r) {
    if (node == -1) return INF;
    int res = tree[node].line.eval(x);
    if (l == r) return res;
    int mid = (l + r) / 2;
    if (x <= mid) return min(res, query(x, tree[node].lc, l, mid));
    return min(res, query(x, tree[node].rc, mid + 1, r));
  }
};
\end{lstlisting}



\catsection{9. Greedy, Searching \& Techniques}

\algotitle{9.1 Binary Search (First True \& Last True)}
\algodesc{Standard monotonic binary search implementations preventing infinite loops.}
\begin{lstlisting}
// First True (Smallest x where check(x) == true):
int low = 0, high = 1e18, ans = -1;
while (low <= high) {
  int mid = low + (high - low) / 2;
  if (check(mid)) { ans = mid; high = mid - 1; }
  else low = mid + 1;
}

// Last True (Largest x where check(x) == true):
low = 0, high = 1e18, ans = -1;
while (low <= high) {
  int mid = low + (high - low) / 2;
  if (check(mid)) { ans = mid; low = mid + 1; }
  else high = mid - 1;
}
\end{lstlisting}

\algotitle{9.2 Floating Point Binary Search}
\algodesc{Fixing iteration count (e.g. 100 steps) guarantees precision $2^{-100} < 10^{-30}$ without precision loops.}
\begin{lstlisting}
double low = 0.0, high = 1e9;
for (int iter = 0; iter < 100; iter++) {
  double mid = (low + high) / 2.0;
  if (check(mid)) high = mid;
  else low = mid;
}
\end{lstlisting}

\algotitle{9.3 Ternary Search (Unimodal Functions)}
\algodesc{Finds maximum/minimum of a strictly convex/concave function in $O(\log \frac{\text{range}}{\epsilon})$.}
\begin{lstlisting}
// Continuous Ternary Search (Finding minimum):
double l = -1e9, r = 1e9;
for (int i = 0; i < 200; i++) {
  double m1 = l + (r - l) / 3.0;
  double m2 = r - (r - l) / 3.0;
  if (f(m1) > f(m2)) l = m1;
  else r = m2;
}

// Discrete Ternary Search (Integers):
int low = 1, high = 1e9;
while (high - low > 2) {
  int m1 = low + (high - low) / 3;
  int m2 = high - (high - low) / 3;
  if (f(m1) < f(m2)) high = m2;
  else low = m1;
}
int ans = f(low);
for (int i = low + 1; i <= high; i++) ans = min(ans, f(i));
\end{lstlisting}

\algotitle{9.4 Two Pointers Technique}
\algodesc{Maintains linear relationships across two indices moving in the same or opposite directions.}
\begin{lstlisting}
// Two Sum in Sorted Array:
int l = 0, r = n - 1;
while (l < r) {
  int sum = a[l] + a[r];
  if (sum == target) { /* Found */ break; }
  else if (sum < target) l++;
  else r--;
}
\end{lstlisting}

\algotitle{9.5 Sliding Window (AtMost Trick)}
\algodesc{Solves exact count queries: $\text{Count}(K) = \text{AtMost}(K) - \text{AtMost}(K - 1)$.}
\begin{lstlisting}
// Number of subarrays with at most K distinct elements:
int at_most_k_distinct(const vector<int>& a, int k) {
  if (k < 0) return 0;
  unordered_map<int, int> freq;
  int l = 0, ans = 0;
  for (int r = 0; r < (int)a.size(); r++) {
    if (freq[a[r]]++ == 0) k--;
    while (k < 0) {
      if (--freq[a[l]] == 0) k++;
      l++;
    }
    ans += (r - l + 1);
  }
  return ans;
}
int exact_k_distinct(const vector<int>& a, int k) {
  return at_most_k_distinct(a, k) - at_most_k_distinct(a, k - 1);
}
\end{lstlisting}

\algotitle{9.6 Meet in the Middle}
\algodesc{Splits $N \le 40$ search spaces into two halves of size $N/2$, then combines via sorting/binary search.}
\begin{lstlisting}
void generate_subset_sums(const vector<int>& a, vector<int>& sums) {
  int n = a.size();
  for (int mask = 0; mask < (1 << n); mask++) {
    int cur = 0;
    for (int i = 0; i < n; i++) if (mask & (1 << i)) cur += a[i];
    sums.push_back(cur);
  }
  sort(sums.begin(), sums.end());
}
int meet_in_the_middle(const vector<int>& a, int target) {
  int n = a.size();
  vector<int> left(a.begin(), a.begin() + n / 2);
  vector<int> right(a.begin() + n / 2, a.end());
  vector<int> sumL, sumR;
  generate_subset_sums(left, sumL);
  generate_subset_sums(right, sumR);
  int best = 0;
  for (int s : sumL) {
    if (s <= target) {
      auto it = upper_bound(sumR.begin(), sumR.end(), target - s);
      if (it != sumR.begin()) best = max(best, s + *(--it));
    }
  }
  return best;
}
\end{lstlisting}



\catsection{10. Computational Geometry}

\algotitle{10.1 Point Structure \& Vector Fundamentals}
\algodesc{2D point primitive with dot product, cross product, orientation, and distance.}
\begin{lstlisting}
const double EPS = 1e-9;
const double PI = acos(-1.0);
struct Point {
  double x, y;
  Point(double x = 0, double y = 0) : x(x), y(y) {}
  Point operator+(const Point& p) const { return Point(x + p.x, y + p.y); }
  Point operator-(const Point& p) const { return Point(x - p.x, y - p.y); }
  Point operator*(double k) const { return Point(x * k, y * k); }
  Point operator/(double k) const { return Point(x / k, y / k); }
  bool operator<(const Point& p) const {
    if (abs(x - p.x) > EPS) return x < p.x;
    return y < p.y - EPS;
  }
};
double dot(Point a, Point b) { return a.x * b.x + a.y * b.y; }
double cross(Point a, Point b) { return a.x * b.y - a.y * b.x; }
double dist_sq(Point a, Point b) { return (a.x - b.x)*(a.x - b.x) + (a.y - b.y)*(a.y - b.y); }
double dist(Point a, Point b) { return hypot(a.x - b.x, a.y - b.y); }
// Orientation: +1 = CCW, -1 = CW, 0 = Collinear
int orient(Point a, Point b, Point c) {
  double v = cross(b - a, c - a);
  if (v > EPS) return 1;
  if (v < -EPS) return -1;
  return 0;
}
Point rotate_point(Point p, double angle) {
  return Point(p.x * cos(angle) - p.y * sin(angle), p.x * sin(angle) + p.y * cos(angle));
}
\end{lstlisting}

\algotitle{10.2 Lines, Segments \& Intersections}
\algodesc{Segment intersection and point-to-segment distance.}
\begin{lstlisting}
bool on_segment(Point p, Point a, Point b) {
  return orient(a, b, p) == 0 &&
         min(a.x, b.x) <= p.x + EPS && p.x <= max(a.x, b.x) + EPS &&
         min(a.y, b.y) <= p.y + EPS && p.y <= max(a.y, b.y) + EPS;
}
bool intersect_segments(Point a, Point b, Point c, Point d) {
  int o1 = orient(a, b, c), o2 = orient(a, b, d);
  int o3 = orient(c, d, a), o4 = orient(c, d, b);
  if (o1 != o2 && o3 != o4) return true;
  if (on_segment(c, a, b) || on_segment(d, a, b)) return true;
  if (on_segment(a, c, d) || on_segment(b, c, d)) return true;
  return false;
}
Point line_intersection(Point a, Point b, Point c, Point d) {
  double cp = cross(b - a, d - c);
  return a + (b - a) * (cross(c - a, d - c) / cp);
}
double dist_to_segment(Point p, Point a, Point b) {
  if (dot(b - a, p - a) <= 0) return dist(p, a);
  if (dot(a - b, p - b) <= 0) return dist(p, b);
  return abs(cross(b - a, p - a)) / dist(a, b);
}
\end{lstlisting}

\algotitle{10.3 Polygon Area (Shoelace Formula)}
\algodesc{Computes oriented area of polygon in $O(N)$. Vertices must be ordered CW or CCW.}
\begin{lstlisting}
double polygon_area(const vector<Point>& p) {
  double area = 0;
  int n = p.size();
  for (int i = 0; i < n; i++) area += cross(p[i], p[(i + 1) % n]);
  return abs(area) / 2.0;
}
\end{lstlisting}

\algotitle{10.4 Convex Hull (Monotonic Chain / Andrew's)}
\algodesc{Finds vertices of the convex hull in $O(N \log N)$ counterclockwise order.}
\begin{lstlisting}
vector<Point> convex_hull(vector<Point> pts) {
  int n = pts.size(), k = 0;
  if (n <= 2) return pts;
  vector<Point> h(2 * n);
  sort(pts.begin(), pts.end());
  for (int i = 0; i < n; h[k++] = pts[i++]) {
    while (k >= 2 && cross(h[k - 1] - h[k - 2], pts[i] - h[k - 2]) <= 0) k--;
  }
  for (int i = n - 2, t = k + 1; i >= 0; h[k++] = pts[i--]) {
    while (k >= t && cross(h[k - 1] - h[k - 2], pts[i] - h[k - 2]) <= 0) k--;
  }
  h.resize(k - 1);
  return h;
}
\end{lstlisting}

\algotitle{10.5 Rotating Calipers (Polygon Diameter)}
\algodesc{Finds the maximum distance between any two points of a convex polygon in $O(N)$.}
\begin{lstlisting}
double rotating_calipers_diameter(const vector<Point>& h) {
  int n = h.size();
  if (n <= 1) return 0;
  if (n == 2) return dist(h[0], h[1]);
  double max_d = 0;
  for (int i = 0, j = 1; i < n; i++) {
    while (cross(h[(i + 1) % n] - h[i], h[(j + 1) % n] - h[j]) > 0) j = (j + 1) % n;
    max_d = max({max_d, dist(h[i], h[j]), dist(h[(i + 1) % n], h[j])});
  }
  return max_d;
}
\end{lstlisting}

\algotitle{10.6 Point in Polygon Checks}
\algodesc{Winding number for general polygon $O(N)$; binary search for convex polygon $O(\log N)$.}
\begin{lstlisting}
// General Polygon: 0 = outside, 1 = strictly inside, 2 = on boundary
int point_in_polygon(const vector<Point>& p, Point q) {
  int n = p.size(); bool inside = false;
  for (int i = 0; i < n; i++) {
    Point a = p[i], b = p[(i + 1) % n];
    if (on_segment(q, a, b)) return 2;
    if ((a.y > q.y) != (b.y > q.y)) {
      double x_cross = a.x + (b.x - a.x) * (q.y - a.y) / (b.y - a.y);
      if (q.x < x_cross) inside = !inside;
    }
  }
  return inside ? 1 : 0;
}
\end{lstlisting}

\algotitle{10.7 Minimum Enclosing Circle (Welzl's Algorithm)}
\algodesc{Finds smallest radius circle enclosing all points in randomized expected $O(N)$.}
\begin{lstlisting}
struct Circle { Point c; double r; };
Circle circumcircle_3pts(Point a, Point b, Point c) {
  Point d1 = b - a, d2 = c - a;
  double d = 2.0 * cross(d1, d2);
  Point center(a.x + (d1.x*d1.x + d1.y*d1.y)*d2.y / d - (d2.x*d2.x + d2.y*d2.y)*d1.y / d,
               a.y + (d2.x*d2.x + d2.y*d2.y)*d1.x / d - (d1.x*d1.x + d1.y*d1.y)*d2.x / d);
  return {center, dist(center, a)};
}
Circle welzl(vector<Point>& pts, vector<Point> r, int n) {
  if (n == 0 || r.size() == 3) {
    if (r.empty()) return {{0, 0}, 0};
    if (r.size() == 1) return {r[0], 0};
    if (r.size() == 2) return {(r[0] + r[1]) / 2.0, dist(r[0], r[1]) / 2.0};
    return circumcircle_3pts(r[0], r[1], r[2]);
  }
  Point p = pts[n - 1];
  Circle c = welzl(pts, r, n - 1);
  if (dist(c.c, p) <= c.r + EPS) return c;
  r.push_back(p);
  return welzl(pts, r, n - 1);
}
Circle min_enclosing_circle(vector<Point> pts) {
  random_shuffle(pts.begin(), pts.end());
  return welzl(pts, {}, pts.size());
}
\end{lstlisting}

\algotitle{10.8 Classical Geometry Theorems \& Formulas}
\algodesc{Pick's theorem, cyclic quadrilaterals, incircle, and circumcircle.}
\begin{itemize}[leftmargin=*,noitemsep,topsep=1pt]
  \item Pick's Theorem: $Area = I + \frac{B}{2} - 1$ \ ($I$: interior lattice points, $B$: boundary points)
  \item Heron's Formula: $Area = \sqrt{s(s - a)(s - b)(s - c)}$ where $s = \frac{a + b + c}{2}$
  \item Triangle Inradius \& Circumradius: $r = \frac{Area}{s}, \quad R = \frac{a \cdot b \cdot c}{4 \cdot Area}$
  \item Brahmagupta's Formula (Cyclic Quad): $Area = \sqrt{(s-a)(s-b)(s-c)(s-d)}$ where $s = \frac{a+b+c+d}{2}$
  \item Fermat Point: Minimizes distance sum to triangle vertices. If all angles $< 120^\circ$, lines to vertices meet at $120^\circ$.
  \item Pythagorean Triples: $a = 2st, \ b = s^2 - t^2, \ c = s^2 + t^2$ with $s > t > 0, \gcd(s, t) = 1, s \not\equiv t \pmod 2$.
\end{itemize}



\catsection{11. Bit Hacks, Game Theory \& STL Reference}

\algotitle{11.1 Bit Manipulation Cheatsheet}
\algodesc{Common bit hacks and GCC builtins for fast binary arithmetic.}
\begin{lstlisting}
// Basic Bit Hacks:
x & -x                   // Lowest set bit (LSB) value
x & (x - 1)              // Clears the lowest set bit
x | (x - 1)              // Sets all trailing zero bits to 1
~x & (x + 1)             // Isolates lowest 0-bit
x && !(x & (x - 1))      // Checks if x is a power of 2
(x >> i) & 1             // Gets the i-th bit (0-indexed)
x |= (1LL << i)          // Sets the i-th bit
x &= ~(1LL << i)         // Clears the i-th bit
x ^= (1LL << i)          // Toggles the i-th bit
(1LL << k) - 1           // Mask of k lowest bits

// GCC Builtins (Append 'll' for 64-bit integers):
__builtin_popcountll(x)  // Number of 1-bits
__builtin_clzll(x)       // Leading zeros (Undefined for x = 0!)
__builtin_ctzll(x)       // Trailing zeros (Undefined for x = 0!)
__builtin_parityll(x)    // Parity of set bits (popcount % 2)
63 - __builtin_clzll(x)  // MSB position (0-indexed)

// Iterate all submasks of mask m (excluding 0):
for (int s = m; s > 0; s = (s - 1) & m) { /* process s */ }

// Next number with the same popcount (GOSPER'S HACK):
int next_popcount(int x) {
  int c = x & -x, r = x + c;
  return r | (((r ^ x) >> 2) / c);
}
\end{lstlisting}

\algotitle{11.2 Range XOR Sum in $O(1)$}
\algodesc{Computes $\bigoplus_{i=0}^n i$ and $\bigoplus_{i=l}^r i$ in $O(1)$.}
\begin{lstlisting}
int xor_0_to_n(int n) {
  if (n % 4 == 0) return n;
  if (n % 4 == 1) return 1;
  if (n % 4 == 2) return n + 1;
  return 0;
}
int range_xor(int l, int r) {
  return xor_0_to_n(r) ^ xor_0_to_n(l - 1);
}
\end{lstlisting}

\algotitle{11.3 Zobrist Hashing (XOR Hashing)}
\algodesc{Assigns uniform random 64-bit integers to distinct values; hash of multiset is XOR sum of elements.}
\begin{itemize}[leftmargin=*,noitemsep,topsep=1pt]
  \item Subarray Permutation / Anagram Check: Two ranges $[l_1, r_1]$ and $[l_2, r_2]$ have identical multiset elements iff $\bigoplus_{i=l_1}^{r_1} \text{rnd}[a[i]] == \bigoplus_{j=l_2}^{r_2} \text{rnd}[b[j]]$.
  \item Tree Isomorphism: Assign random value to depth or subtree sizes, hash by XORing child subtree hashes.
\end{itemize}
\begin{lstlisting}
mt19937_64 rng(chrono::steady_clock::now().time_since_epoch().count());
uint64_t random_val() { return rng(); }
\end{lstlisting}

\algotitle{11.4 Game Theory \& Nim Rules}
\algodesc{Impartial games under the normal play convention (last player to move wins).}
\begin{itemize}[leftmargin=*,noitemsep,topsep=1pt]
  \item Standard Nim: Pile sizes $a_1, a_2, \dots, a_n$. First player wins iff XOR sum $X = a_1 \oplus a_2 \oplus \dots \oplus a_n \ne 0$.
  \item Subtraction Game: Can remove at most $k$ stones from a pile. Equivalent to Nim with pile sizes $a_i \bmod (k + 1)$.
  \item Index-$k$ Nim: Can remove from at most $k$ piles in a single turn. Write pile sizes in binary; sum each bit column modulo $k + 1$. First player wins iff at least one bit column sum $\not\equiv 0 \pmod{k + 1}$.
  \item Misère Nim: Normal play, but the last player to move LOSES.
    \begin{itemize}[noitemsep]
      \item If all piles have size 1: First player wins iff total number of piles is EVEN.
      \item If at least one pile has size $> 1$: Play identically to Standard Nim ($X \ne 0 \implies$ First player wins).
    \end{itemize}
  \item Wythoff's Game: Two piles of size $(n, m)$. Can take any from one, or equal from both. Losing states are $(\lfloor k \phi \rfloor, \lfloor k \phi^2 \rfloor)$ where $\phi = \frac{1 + \sqrt{5}}{2} \approx 1.6180339887$.
\end{itemize}

\algotitle{11.5 Sprague-Grundy Theorem}
\algodesc{Every impartial game position is equivalent to a single Nim pile of size $g(u) = \operatorname{mex}(\{g(v) : u \to v\})$.}
\begin{lstlisting}
int memo_g[N];
int grundy(int state, const vector<int>& valid_moves) {
  if (state == 0) return 0;
  if (memo_g[state] != -1) return memo_g[state];
  unordered_set<int> transitions;
  for (int m : valid_moves) {
    if (state >= m) transitions.insert(grundy(state - m, valid_moves));
  }
  int mex = 0;
  while (transitions.count(mex)) mex++;
  return memo_g[state] = mex;
}
\end{lstlisting}

\algotitle{11.6 STL Container Iteration \& Erase}
\algodesc{Safely deleting elements while iterating across std::map and std::set without invalidating iterators.}
\begin{lstlisting}
for (auto it = s.begin(); it != s.end(); ) {
  if (condition(*it)) it = s.erase(it); // Erase returns next valid iterator
  else ++it;
}
\end{lstlisting}

\algotitle{11.7 Direction Arrays (Grid Problems)}
\algodesc{Standard offset deltas for 4-directional, 8-directional, and Knight moves on 2D grids.}
\begin{lstlisting}
// 4-Directional:
const int dx4[4] = {-1, 0, 1, 0};
const int dy4[4] = {0, 1, 0, -1};

// 8-Directional:
const int dx8[8] = {-1, -1, -1, 0, 0, 1, 1, 1};
const int dy8[8] = {-1, 0, 1, -1, 1, -1, 0, 1};

// Knight's Moves:
const int kx[8] = {-2, -2, -1, -1, 1, 1, 2, 2};
const int ky[8] = {-1, 1, -2, 2, -2, 2, -1, 1};
\end{lstlisting}

\algotitle{11.8 "When to Use What" Contest Cheat Sheet}
\algodesc{Quick problem pattern recognition table from IUPC notebook.}
\begin{tabular}{|l|l|}
\hline
\textbf{Problem Pattern} & \textbf{Algorithm / DS} \\
\hline
Range sum with point updates & Fenwick Tree (BIT) / SegTree \\
Range add, Range sum/min & Lazy Segment Tree \\
Range minimum query (Static) & Sparse Table $O(1)$ \\
Count pairs with XOR condition & Binary Trie / Basis Vector \\
Offline range distinct/frequency & Mo's Algorithm with snake sort \\
Shortest path (unweighted) & BFS \\
Shortest path ($0/1$ weights) & 0-1 BFS with Deque \\
Shortest path (positive weights) & Dijkstra with priority\_queue \\
Shortest path (negative edges) & Bellman-Ford / SPFA \\
All-pairs shortest paths & Floyd-Warshall $O(V^3)$ \\
Connected components in directed & Kosaraju / Tarjan SCC \\
Tree path sum/update & Heavy-Light Decomposition (HLD) \\
Lowest Common Ancestor & Binary Lifting $O(\log N)$ \\
Offline subtree queries & DSU on Tree (Sack) \\
Pattern matching in text & KMP / Z-Algorithm / Hashing \\
Longest palindrome substring & Manacher's Algorithm $O(N)$ \\
Divisors / Prime factorization & Linear Sieve SPF $O(N)$ \\
$a^b \pmod M$ & Binary Exponentiation $O(\log b)$ \\
Linear recurrence $k$-th term & Matrix Exponentiation $O(N^3 \log k)$ \\
Optimal split interval DP & Knuth Optimization $O(N^2)$ \\
Dynamic line min/max queries & Li Chao Tree $O(\log C)$ \\
\hline
\end{tabular}


\end{multicols*}
\end{document}
